\documentclass[11pt]{article}

\usepackage[letterpaper,margin=0.82in,headheight=15pt]{geometry}
\usepackage[T1]{fontenc}
\usepackage{lmodern}
\usepackage{microtype}
\usepackage{amsmath,amssymb,amsthm,bm,mathtools}
\usepackage{booktabs,longtable,tabularx,array,multirow}
\usepackage{enumitem}
\usepackage{xcolor}
\usepackage{fancyhdr}
\usepackage{graphicx}
\usepackage{tikz-cd}
\usepackage{hyperref}
\usepackage[nameinlink,noabbrev]{cleveref}

\definecolor{navy}{HTML}{17324D}
\definecolor{teal}{HTML}{147D92}
\definecolor{orange}{HTML}{D47A27}
\definecolor{softblue}{HTML}{EEF4F7}
\definecolor{softgray}{HTML}{F3F5F6}
\definecolor{darkgray}{HTML}{46545E}
\definecolor{softorange}{HTML}{FFF4E8}

\hypersetup{
  colorlinks=true,
  linkcolor=navy,
  citecolor=teal,
  urlcolor=teal,
  pdftitle={Volume V: Regulator-to-QCD Matching, Rank-Aware Evolution, and Process Factorization},
  pdfauthor={DeuteronWigner project}
}

\pagestyle{fancy}
\fancyhf{}
\fancyhead[L]{\small\color{darkgray}DeuteronWigner formalism - Volume V}
\fancyhead[R]{\small\color{darkgray}Matching, evolution, and process factorization}
\fancyfoot[L]{\small\color{darkgray}30 July 2026}
\fancyfoot[R]{\small\color{darkgray}\thepage}
\renewcommand{\headrulewidth}{0.4pt}

\setlist[itemize]{leftmargin=1.5em,itemsep=2pt,topsep=3pt}
\setlist[enumerate]{leftmargin=1.7em,itemsep=2pt,topsep=3pt}
\setlength{\parindent}{0pt}
\setlength{\parskip}{5pt}
\setlength{\emergencystretch}{2em}
\renewcommand{\arraystretch}{1.16}
\allowdisplaybreaks

\newtheorem{definition}{Definition}[section]
\newtheorem{axiom}[definition]{Axiom}
\newtheorem{principle}[definition]{Design principle}
\newtheorem{requirement}[definition]{Requirement}
\newtheorem{proposition}[definition]{Proposition}
\newtheorem{theorem}[definition]{Theorem}
\newtheorem{lemma}[definition]{Lemma}
\newtheorem{corollary}[definition]{Corollary}
\newtheorem{remark}[definition]{Remark}

\crefname{definition}{definition}{definitions}
\crefname{axiom}{axiom}{axioms}
\crefname{principle}{principle}{principles}
\crefname{requirement}{requirement}{requirements}
\crefname{proposition}{proposition}{propositions}
\crefname{theorem}{theorem}{theorems}
\crefname{lemma}{lemma}{lemmas}
\crefname{corollary}{corollary}{corollaries}
\crefname{remark}{remark}{remarks}

\newcommand{\kT}{\bm{k}_{T}}
\newcommand{\ellT}{\bm{\ell}_{T}}
\newcommand{\qT}{\bm{q}_{T}}
\newcommand{\PT}{\bm{P}_{T}}
\newcommand{\DT}{\bm{\Delta}_{T}}
\newcommand{\bD}{\bm{b}_{\Delta}}
\newcommand{\bM}{\bm{b}_{\mathrm{TMD}}}
\newcommand{\dd}{\mathrm{d}}
\newcommand{\Tr}{\operatorname{Tr}}
\newcommand{\tr}{\operatorname{tr}}
\newcommand{\Span}{\operatorname{span}}
\newcommand{\End}{\operatorname{End}}
\newcommand{\Hom}{\operatorname{Hom}}
\newcommand{\Id}{\operatorname{Id}}
\newcommand{\Ker}{\operatorname{Ker}}
\newcommand{\Ran}{\operatorname{Ran}}
\newcommand{\rank}{\operatorname{rank}}
\newcommand{\diag}{\operatorname{diag}}
\newcommand{\Hil}{\mathcal H}
\newcommand{\Alg}{\mathcal A}
\newcommand{\Rep}{\mathcal V}
\newcommand{\LF}{\mathrm{LF}}
\newcommand{\TMD}{\mathrm{TMD}}
\newcommand{\QCD}{\mathrm{QCD}}
\newcommand{\GTMD}{\mathrm{GTMD}}
\newcommand{\GPD}{\mathrm{GPD}}
\newcommand{\EFT}{\mathrm{EFT}}
\newcommand{\Amp}{\mathbf{Amp}}
\newcommand{\Dens}{\mathbf{Dens}}
\newcommand{\Match}{\mathbf{Match}}
\newcommand{\Red}{\mathbf{Red}}
\newcommand{\Proc}{\mathbf{Proc}}
\newcommand{\cO}{\mathcal O}
\newcommand{\cM}{\mathcal M}
\newcommand{\cR}{\mathcal R}
\newcommand{\cS}{\mathcal S}
\newcommand{\cV}{\mathcal V}
\newcommand{\cW}{\mathcal W}
\newcommand{\cE}{\mathcal E}
\newcommand{\cG}{\mathcal G}
\newcommand{\cK}{\mathcal K}
\newcommand{\cQ}{\mathcal Q}
\newcommand{\cP}{\mathcal P}
\newcommand{\cC}{\mathcal C}
\newcommand{\cZ}{\mathcal Z}
\newcommand{\cD}{\mathcal D}
\newcommand{\cN}{\mathcal N}
\newcommand{\cU}{\mathcal U}
\newcommand{\cF}{\mathcal F}
\newcommand{\cA}{\mathcal A}
\newcommand{\cB}{\mathcal B}
\newcommand{\cL}{\mathcal L}
\newcommand{\cJ}{\mathcal J}
\newcommand{\cI}{\mathcal I}
\newcommand{\eps}{\varepsilon}
\newcommand{\bbone}{\bm 1}
\newcommand{\abs}[1]{\left|#1\right|}
\newcommand{\norm}[1]{\left\lVert#1\right\rVert}
\newcommand{\inner}[2]{\left\langle #1\middle|#2\right\rangle}
\newcommand{\avg}[1]{\left\langle #1\right\rangle}
\newcommand{\comm}[2]{\left[#1,#2\right]}
\newcommand{\acom}[2]{\left\{#1,#2\right\}}
\newcommand{\otimesx}{\mathop{\otimes}_{x}}
\newcommand{\MSbar}{\overline{\mathrm{MS}}}
\newcommand{\ren}{\mathrm{ren}}
\newcommand{\unsub}{\mathrm{unsub}}
\newcommand{\soft}{\mathrm{soft}}
\newcommand{\rap}{\mathrm{rap}}
\newcommand{\UV}{\mathrm{UV}}
\newcommand{\IR}{\mathrm{IR}}
\newcommand{\pert}{\mathrm{pert}}
\newcommand{\NP}{\mathrm{NP}}
\newcommand{\FO}{\mathrm{FO}}
\newcommand{\asy}{\mathrm{asy}}
\newcommand{\evol}{\mathrm{evol}}
\newcommand{\proc}{\mathrm{proc}}
\newcommand{\tw}{\mathrm{tw}}
\newcommand{\NS}{\mathrm{NS}}
\newcommand{\sing}{\mathrm{S}}
\newcommand{\CS}{\mathrm{CS}}
\newcommand{\DY}{\mathrm{DY}}
\newcommand{\SIDIS}{\mathrm{SIDIS}}
\newcommand{\EMT}{\mathrm{EMT}}
\newcommand{\IA}{\mathrm{IA}}
\newcommand{\pT}{\bm{p}_{T}}
\newcommand{\shortd}{\mathrm{short}}
\newcommand{\nuc}{\mathrm{nuc}}
\newcommand{\even}{\mathrm{even}}
\newcommand{\odd}{\mathrm{odd}}

\newcolumntype{Y}{>{\raggedright\arraybackslash}X}
\newcolumntype{L}[1]{>{\raggedright\arraybackslash}p{#1}}

\title{\vspace{-1.0cm}
{\color{navy}\bfseries Volume V: Regulator-to-QCD Matching,\\[2mm]
Rank-Aware Evolution, and Process Factorization}\\[4mm]
\large Operator matching, two-scale renormalization, coupled collinear evolution,\\
process-qualified TMD observables, and auditable \(W+Y\) assembly}
\author{DeuteronWigner project}
\date{30 July 2026}

\begin{document}
\maketitle

\begin{center}
\begin{minipage}{0.94\textwidth}
\colorbox{softblue}{\parbox{0.96\textwidth}{
\textbf{Status and purpose.}
Volumes~I--IV define the regulated light-front state spaces, common nucleon
GTMD overlaps, dynamical Wilson-line phases, and the matched spin-1 nuclear
state.  Those objects are not yet physical TMDs in a declared continuum QCD
scheme.  The present volume supplies the missing map.  It separates
(i) matching a regulated Hamiltonian operator to a renormalized QCD operator,
(ii) ultraviolet and rapidity evolution in the two-dimensional
\((\mu,\zeta)\) scale plane, and (iii) process-level assembly with hard,
fragmentation, color, measurement, and fixed-order remainder kernels.  A
single rank-aware route must carry every flavor, target polarization,
Wilson-link/color class, nuclear component, and correlated microscopic member.
The accepted fixed-scale canonical model remains a comparison and regression
boundary.  Evolving that boundary is useful, but it does not by itself turn
heterogeneous fitted inputs into one microscopic prediction.
}}
\end{minipage}
\end{center}

\begin{abstract}
A low-resolution light-front Hamiltonian wave function is not, by itself, a
transverse-momentum-dependent distribution in a specified QCD factorization
scheme.  Its bilocal operators carry regulator, model-space, and truncation
dependence; a TMD additionally requires ultraviolet renormalization, rapidity
renormalization, soft subtraction, a Wilson-path identity, and two evolution
scales.  This volume constructs the regulator-to-QCD matching and evolution
layer required to turn the common microscopic quark, antiquark, and gluon
parents of Volumes~II--IV into scheme-qualified distributions and
process-level observables.

The matching is formulated as a typed linear map between operator bases,
not as one free normalization per named TMD.  At a reference scale it relates
regulated light-front matrix elements to renormalized QCD GTMD, TMD, GPD,
current, and energy--momentum operators, including operator mixing, soft and
rapidity factors, power corrections, and explicit Hamiltonian truncation
errors.  When the Hamiltonian cutoff is not perturbatively high, the matching
coefficient is a controlled combination of perturbation theory, step scaling,
local-current constraints, lattice or continuum matrix elements, and a small
shared discrepancy basis.  The same matching parameter must propagate to
every projection of the operator on which it acts.

We define the two-scale TMD evolution equations using a convention in which
the mixed derivatives are exactly compatible, introduce the evolution
one-form and its finite-order curl defect, and distinguish exact path
independence from the ambiguity created by truncated perturbation theory.  A
scheme-neutral formulation accommodates direct contour evolution,
\(\zeta\)-prescription evolution, and profile-based implementations while
requiring an explicit accuracy manifest.  The Collins--Soper kernel is
separated into perturbative and nonperturbative domains; target and spin
polarization reside in boundary matrix elements and small-\(b\) matching,
not in an independently fitted evolution kernel for every TMD.  Quark and
gluon kernels remain separate color-representation objects, and generalized
Casimir relations are used only to their verified perturbative order.

The small-\(b_{\mathrm{TMD}}\) operator product expansion is organized by
parton species, operator polarization, transverse rank, target channel,
Wilson-link class, and collinear twist.  Twist-2 quark, antiquark, and gluon
channels; twist-3 Sivers, Boer--Mulders, and worm-gear channels; linearly
polarized gluons; tensor-polarized spin-1 distributions; and independent
gluon \(f\)- and \(d\)-type sectors are assigned to explicit collinear
operator bases.  Nonsinglet, singlet quark--gluon, helicity, transversity,
double-helicity-flip, and multiparton twist-3 evolution are kept distinct.
Heavy-flavor thresholds are matching maps, not silent changes of flavor
count.

At the process layer, Drell--Yan, SIDIS, inclusive DIS, and gluon-sensitive
channels receive typed factorization records containing the hard kernel,
link map, color weights, fragmentation or partner TMD, measurement function,
Glauber status, and factorization status.  The observable is assembled as
\(W+Y\), with
\(Y=\sigma_{\FO}-[W]_{\asy,\FO}\) at the same perturbative order and in the
same scheme.  Rank-resolved Fourier--Bessel kernels are used for every angular
harmonic, so a scalar \(J_0\) transform can never be applied to a nonzero-rank
coefficient.  We provide formal commutation results for evolution and linear
nuclear convolution, scheme-covariance relations, uncertainty and convergence
requirements, software interfaces, analytic benchmarks, and fail-closed
acceptance gates.  The result is a reproducible bridge from the regulated
microscopic state to multi-scale QCD observables without erasing flavor,
spin, rank, link/color, nuclear, or member identity.
\end{abstract}

\tableofcontents

\section{Scope, interfaces, and status language}
\label{sec:scope}

\subsection{Input and output contract}

The input to this volume is not a table of named TMDs.  It is a versioned
family of regulated operator matrix elements,
\begin{equation}
 \widetilde{\bm W}^{\LF}_{h,r}
 =
 \left\{
 \langle P',\Lambda'|
 \widetilde{\cO}^{\LF}_{A,r}(x,\bM,\DT)
 |P,\Lambda\rangle
 \right\}_{A},
 \label{eq:lf-parent-vector}
\end{equation}
where \(r\) records the Hamiltonian regulator, basis, Fock truncation,
renormalization conditions, and microscopic member, and \(A\) runs over a
closed operator basis.  For a deuteron, the matrix element is evaluated in the
matched nuclear state of Volume~IV.  The impact variable in
\cref{eq:lf-parent-vector} is \(\bM\), Fourier conjugate to \(\kT\); the
partonic imaging coordinate \(\bD\), conjugate to \(\DT\), remains a distinct
type.

The output is a process-qualified observable or a renormalized intermediate
object,
\begin{equation}
 \widetilde{\bm W}^{\QCD}_{h}
 (\mu,\zeta;\mathfrak s),
 \qquad
 \widetilde{\bm F}_{h}
 (\mu,\zeta;\mathfrak s),
 \qquad
 \frac{\dd\sigma_{\proc}}{\dd\cP}
 =W_{\proc}+Y_{\proc}+\delta_{\mathrm{pow}},
 \label{eq:volume5-outputs}
\end{equation}
with declared QCD scheme \(\mathfrak s\), scales, process map, perturbative
accuracy, and power-correction status.

\begin{principle}[Three transformations, three meanings]
The regulator-to-QCD matching map, TMD scale evolution, and process
factorization map are different morphism classes.  Matching changes the
operator description, evolution changes scales within one description, and
process factorization contracts universal or process-qualified operator
matrix elements with short-distance and measurement kernels.  No software
callable may silently serve all three roles.
\end{principle}

\subsection{What this volume does not claim}

This volume does not prove a factorization theorem for every proposed
gluon-sensitive process, compute every small-\(b\) coefficient, or determine
the nonperturbative Collins--Soper kernel from first principles.  It defines
the exact interfaces at which those ingredients enter and the conditions
under which a result may be called authoritative.  In particular:
\begin{itemize}
\item a formal Fourier transform is not a renormalization prescription;
\item a smooth evolution curve is not evidence that matching is correct;
\item a common value of \(Q\) does not imply a common TMD scheme;
\item a successful \(W\) term at small \(q_T\) does not determine the
      fixed-order \(Y\) term;
\item a process label does not determine a gluon \(f/d\) mixture without a
      hard-color calculation;
\item pointwise positivity of a regulated overlap is not assumed to survive
      every finite-order subtracted scheme representation.
\end{itemize}

\subsection{Evidence classes}

Each matching or evolution ingredient is assigned one of:
\begin{longtable}{L{0.25\textwidth}L{0.65\textwidth}}
\toprule
Status & Meaning\\
\midrule
\endhead
\texttt{EXACT} & Algebraic identity, symmetry, type rule, or all-order
consistency condition within the declared definitions.\\
\texttt{PERTURBATIVE} & Coefficient or anomalous dimension known to a stated
order in \(\alpha_s\) and in a stated scheme.\\
\texttt{NONPERTURBATIVE\_INPUT} & Common kernel or matrix element constrained
by lattice QCD, phenomenology, or a microscopic calculation.\\
\texttt{MATCHED\_MODEL} & Shared low-dimensional matching or discrepancy
operator calibrated to several matrix elements.\\
\texttt{EXPLORATORY} & Process, coefficient, or power correction used only for
sensitivity studies.\\
\texttt{UNAVAILABLE} & Required ingredient absent; authoritative composition
must fail closed.\\
\bottomrule
\end{longtable}
Passing a numerical gate does not promote an exploratory coefficient to a
perturbative theorem or a fitted kernel to a first-principles result.

\section{Convention, scheme, and operator-identity contract}
\label{sec:types}

\subsection{Transverse coordinates and rank}

The four transverse variables needed by this layer are
\begin{equation}
 \begin{aligned}
 \DT&\leftrightarrow\bD,
 &\qquad
 \kT&\leftrightarrow\bM,\\
 \qT&=\text{measured recoil},
 &\qquad
 \ellT&=\text{loop or emitted momentum}.
 \end{aligned}
 \label{eq:coordinate-contract}
\end{equation}
Equal dimensions do not permit substitution.  The partonic Wigner transform
acts in \((\DT,\bD)\), whereas TMD evolution and process Fourier transforms
act in \((\kT,\bM)\).

A transverse-rank identity is
\begin{equation}
 \mathfrak r=
 \left(
 m,\,r,\,\mathcal T_r,\,n_k,\,n_b,\,M_{\rm ref},\,
 J_{|m|},\,\varphi_{\rm FT}
 \right),
 \label{eq:rank-identity}
\end{equation}
where \(m\) is the \(SO(2)\) harmonic weight, \(r\) the tensor rank,
\(\mathcal T_r\) the symmetric-traceless basis, \(n_k,n_b\) the extracted
powers, \(M_{\rm ref}\) the reference mass, and \(\varphi_{\rm FT}\) the
Fourier phase convention.

\begin{requirement}[Rank is part of operator identity]
Two scalar coefficient arrays with different \(\mathfrak r\) are different
operators even when they share the same numerical grid and dimensions.  An
explicit rank adapter is required to compare them.
\end{requirement}

\subsection{QCD TMD scheme tuple}

A complete TMD scheme identifier is
\begin{equation}
 \mathfrak s_{\TMD}=
 \left(
 \mathfrak s_{\UV},
 \mathcal R_{\rap},
 \mathcal S_{\soft},
 \mathcal C_{\rm staple},
 R_{\rm color},
 \mathfrak n_{\gamma_5},
 \mathfrak n_{\epsilon_T},
 \mathfrak n_{\rm FT}
 \right).
 \label{eq:tmd-scheme-tuple}
\end{equation}
It records the ultraviolet subtraction, rapidity regulator, soft-factor
partition, staple/cusp convention, color representation, \(\gamma_5\)
implementation, transverse Levi-Civita convention, and Fourier normalization.
The project reference realization is compatible with \(\MSbar\) ultraviolet
renormalization, a declared rapidity regulator, and square-root soft
subtraction, but no formula in this volume assumes that a finite factor is
identical across all schemes.

The full operator identity extends the tuple introduced in Volumes~0--III:
\begin{equation}
 \mathfrak o=
 \left(
 a,\Gamma,P_{\rm in},P_{\rm out},\gamma,R_{\rm color},
 \mathfrak g_{f/d},\mathfrak r,
 \mathfrak s_{\TMD},\mu,\zeta,
 r_{\LF},\nu_{\rm nuc},m_{\rm micro}
 \right).
 \label{eq:full-operator-identity}
\end{equation}
Here \(\mathfrak g_{f/d}\) includes the ordered gluon-link pair and independent
\(f^{abc}\) or \(d^{abc}\) contraction where applicable, \(r_{\LF}\) is the
Hamiltonian regulator/truncation label, \(\nu_{\rm nuc}\) is the nuclear
sector or operator block, and \(m_{\rm micro}\) is the common microscopic
member.

\begin{requirement}[No production use of unspecified identity]
An operator may be represented with an explicit \texttt{UNSPECIFIED} field
during migration, but it cannot enter a matching, evolution, or process map
that mathematically depends on that field.  No default scheme, link, rank,
color class, or scale is inferred from a function name.
\end{requirement}

\subsection{Scheme transformations}

Let \(\mathfrak s\) and \(\mathfrak s'\) be two complete schemes for the same
physical operator.  A finite scheme transformation has the form
\begin{equation}
 \widetilde{\bm F}^{\mathfrak s'}
 =\bm Z^{\mathfrak s'\leftarrow\mathfrak s}
  \otimesx
  \widetilde{\bm F}^{\mathfrak s},
 \label{eq:finite-scheme-transform}
\end{equation}
where \(\bm Z\) may be a matrix in operator or color space.  For a factorized
observable with two TMD factors,
\begin{equation}
 \bm H^{\mathfrak s'}
 =
 \bm H^{\mathfrak s}
 \left(\bm Z_1^{\mathfrak s'\leftarrow\mathfrak s}\right)^{-1}
 \left(\bm Z_2^{\mathfrak s'\leftarrow\mathfrak s}\right)^{-1},
 \label{eq:hard-scheme-transform}
\end{equation}
with the appropriate index ordering, so the complete cross section is
invariant up to the declared truncation order.

\begin{proposition}[Scheme covariance]
If the TMDs, hard factor, soft partition, and fixed-order subtraction are all
transformed with the same finite map, then the process observable is unchanged
through the perturbative order to which the transformation is known.
\end{proposition}

\begin{proof}
Insert \cref{eq:finite-scheme-transform,eq:hard-scheme-transform} into the
factorized contraction.  The finite maps cancel pairwise.  Any residual is of
higher perturbative order or belongs to an explicitly omitted operator or
power-correction sector.  The same transformation must be applied to the
asymptotic expansion used in the \(Y\) term; otherwise the cancellation fails.
\end{proof}

\section{From regulated light-front operators to QCD operators}
\label{sec:lf-to-qcd}

\subsection{Three matching stages}

The complete bridge is decomposed as
\begin{equation}
 \cO^{\LF}_{r}
 \xrightarrow{\ \cZ_{\rm state/op}\ }
 \cO^{\rm cont,unsub}
 \xrightarrow{\ \cZ_{\UV}\cR_{\rap}\cS_{\soft}^{-1/2}\ }
 \cO^{\QCD}_{\TMD}(\mu,\zeta;\mathfrak s)
 \xrightarrow{\ \cC_{\rm proc}\ }
 \cO_{\rm measured}.
 \label{eq:three-stage-matching}
\end{equation}
The first arrow removes Hamiltonian regulator and model-space dependence from
the state and operator basis.  The second defines the renormalized,
soft-subtracted QCD TMD or GTMD.  The third is process factorization.  These
arrows may share perturbative ingredients, but they are not interchangeable.

\subsection{General operator-basis matching}

Let \(\bm\cO^{\LF}_r\) and \(\bm\cO^{\QCD}\) be closed bases with identical
conserved quantum numbers.  At reference scales \((\mu_M,\zeta_M)\),
\begin{equation}
 \widetilde{\bm\cO}^{\QCD}
 (\mu_M,\zeta_M)
 =
 \bm Z_{\LF\to\QCD}
 (r;\mu_M,\zeta_M)
 \otimesx
 \widetilde{\bm\cO}^{\LF}_{r}
 +\widetilde{\bm\cO}_{\rm miss}
 +\widetilde{\bm\delta}_{\rm pow}
 +\widetilde{\bm\delta}_{\rm trunc}.
 \label{eq:general-lf-qcd-matching}
\end{equation}
The map can mix quark, gluon, total-derivative, equation-of-motion, and
multiparton operators when allowed by symmetry and twist.  The missing-operator
term is not an arbitrary repair: every component must correspond to a named
operator absent from the retained Hamiltonian or current.

For external calibration states \(|\chi_i\rangle\), matching conditions are
\begin{equation}
 \langle\chi_i|\bm\cO^{\QCD}|\chi_j\rangle
 =
 \bm Z_{\LF\to\QCD}
 \otimesx
 \langle\chi_i|\bm\cO^{\LF}_{r}|\chi_j\rangle
 +\bm\Delta_{ij}(r).
 \label{eq:matrix-element-matching}
\end{equation}
The state set must overconstrain the shared matching parameters.  Matching one
coefficient to one named TMD does not establish an operator map.

\begin{requirement}[Shared matching parameters]
A parameter in \(\bm Z_{\LF\to\QCD}\) belongs to an operator mixing channel,
regulator counterterm, soft/rapidity conversion, or controlled discrepancy
operator.  It propagates to every TMD, GPD, PDF, current, and process
projection of that operator.  One normalization or transverse width per named
TMD is prohibited.
\end{requirement}

\subsection{Perturbative, nonperturbative, and hybrid matching}

If all relevant matching momenta are large compared with \(\Lambda_{\QCD}\),
\(\bm Z_{\LF\to\QCD}\) can be expanded perturbatively.  A low-resolution
Hamiltonian generally does not satisfy this condition over the full \(x\) and
\(\bM\) domain.  The admissible routes are:
\begin{enumerate}
\item \textbf{Perturbative window matching:} determine \(\bm Z\) for
      \(\abs{\bM}^{-1}\), external virtualities, and cutoffs in a common
      perturbative domain.
\item \textbf{Step scaling:} connect a sequence of regulated theories
      \(r_0\prec r_1\prec\cdots\prec r_n\) and match at each overlap window.
\item \textbf{Local-current anchoring:} impose exact charges, Ward identities,
      and renormalized local moments on the same operator family.
\item \textbf{External matrix-element matching:} use lattice, continuum, or
      experimental matrix elements whose scheme conversion is controlled.
\item \textbf{Hybrid matching:} combine the above with a low-dimensional
      discrepancy basis constrained across several channels and regulator
      levels.
\end{enumerate}
A route is part of the model definition and cannot be selected independently
for each observable after seeing the data.

\subsection{Step-scaling system}

For regulator levels \(r\preceq r'\), define
\begin{equation}
 \bm\Sigma_{r'\leftarrow r}
 =
 \bm Z_{\LF\to\QCD}(r')^{-1}
 \bm Z_{\LF\to\QCD}(r),
 \label{eq:step-scaling-map}
\end{equation}
so that matched observables satisfy
\begin{equation}
 \widetilde{\bm W}^{\LF}_{r'}
 =
 \bm\Sigma_{r'\leftarrow r}\otimesx
 \widetilde{\bm W}^{\LF}_{r}
 +\cO(\delta_{r'r}).
 \label{eq:step-scaling-consistency}
\end{equation}
Composition requires
\begin{equation}
 \bm\Sigma_{r''\leftarrow r'}
 \bm\Sigma_{r'\leftarrow r}
 =
 \bm\Sigma_{r''\leftarrow r}
 +\cO(\delta_{r''r'r}).
 \label{eq:step-scaling-cocycle}
\end{equation}
The defect in \cref{eq:step-scaling-cocycle} is a truncation diagnostic, not a
parameter uncertainty.

\subsection{Local limits and conserved operators}

The matching basis must contain the local limits needed for vector, axial,
tensor, and energy--momentum currents.  For example,
\begin{align}
 \int\dd x\,H^q(x,0,t)
 &\longrightarrow
 \langle P'|\bar q\gamma^+q|P\rangle,
 \label{eq:vector-local-limit}\\
 \sum_{a=q,\bar q,g}\int\dd x\,x\,H^a(x,0,t)
 &\longrightarrow
 \langle P'|T^{++}|P\rangle.
 \label{eq:emt-local-limit}
\end{align}
The conserved total vector current and total energy--momentum tensor constrain
mixing matrices and threshold maps.  A matching solution that reproduces a
TMD shape but violates charge or momentum closure is rejected.

\begin{proposition}[No independent local and TMD renormalization]
When a local current is the regulated moment of a bilocal operator, the
counterterms and finite matching factors of the two objects must be related by
the same moment map, up to contact and operator-mixing terms explicitly
present in the basis.
\end{proposition}

\section{Renormalized TMD and GTMD operator definitions}
\label{sec:renormalized-definitions}

\subsection{Unsubtracted, soft, and renormalized objects}

For a quark or gluon representation \(R_a\), let
\(\widetilde W_{a,h}^{\unsub[\gamma]}\) be the regulated bilocal matrix element
and \(S_{R_a}^{[\gamma]}\) the corresponding vacuum soft function.  An abstract
square-root-subtracted definition is
\begin{equation}
 \widetilde W_{a/h}^{\ren[\gamma]}
 (x,\bM,\DT;\mu,\zeta;\mathfrak s)
 =
 Z_{a}^{\UV}(\mu;\mathfrak s)
 R_a^{\rap}(\bM;\mu,\zeta;\mathfrak s)
 \widetilde W_{a/h}^{\unsub[\gamma]}
 \left[S_{R_a}^{[\gamma]}\right]^{-1/2}.
 \label{eq:renormalized-gtmd}
\end{equation}
The exact placement of rapidity and soft factors is scheme dependent.  The
scheme tuple in \cref{eq:tmd-scheme-tuple} makes this dependence explicit.
For color-entangled processes the soft object can be a matrix and a product of
independent single-hadron TMDs may not exist.

The forward TMD is
\begin{equation}
 \widetilde F_{a/h}^{[\gamma]}
 (x,\bM;\mu,\zeta)
 =
 \widetilde W_{a/h}^{[\gamma]}
 (x,\bM,\DT=0;\mu,\zeta).
 \label{eq:forward-tmd-from-gtmd}
\end{equation}
The partonic Wigner transform instead Fourier transforms \(\DT\) into
\(\bD\) and does not alter the rapidity subtraction associated with \(\bM\).

\subsection{GTMD evolution and the GPD reduction}

At leading power, a properly soft-subtracted GTMD has the same ultraviolet and
rapidity evolution structure as the corresponding TMD for a fixed parton
representation; the external transfer and hadron polarization reside in the
matrix element.  However, the formal integral
\begin{equation}
 \int\dd^2\kT\,W_{\TMD\text{-}subtracted}
 \label{eq:naive-gtmd-gpd-integral}
\end{equation}
is not automatically a renormalized GPD.  The \(\bM\to0\) limit contains
ultraviolet singularities, and the staple operator must be matched onto the
appropriate straight-link collinear operator.  The correct square is
\begin{equation}
\begin{tikzcd}[column sep=large,row sep=large]
 \widetilde W_{\LF}^{(r)}
 \arrow[r,"\cZ_{\LF\to\QCD}"]
 \arrow[d,"\mathsf G_{\LF}"']
 &
 \widetilde W_{\QCD}^{\TMD}(\mu,\zeta)
 \arrow[d,"\mathsf OPE_{\bM\to0}"]
 \\
 H_{\LF}^{(r)}
 \arrow[r,"\cZ_{\LF\to\GPD}"']
 &
 H_{\QCD}(\mu)
\end{tikzcd}
 \label{eq:matched-gtmd-gpd-square}
\end{equation}
which must commute within matching, power, truncation, and numerical errors.

\begin{requirement}[No naive marginal promotion]
A regulated \(\kT\) integral may be used as an analytic benchmark, but it is
not labeled a physical GPD, PDF, or current until the matched square
\cref{eq:matched-gtmd-gpd-square} has been supplied.
\end{requirement}

\subsection{Link reversal and evolution}

For process-conjugate future- and past-pointing staples, the link-even and
link-odd boundary conditions satisfy
\begin{equation}
 F_{\even}^{[+]}(Q_0)=F_{\even}^{[-]}(Q_0),
 \qquad
 F_{\odd}^{[+]}(Q_0)=-F_{\odd}^{[-]}(Q_0),
 \label{eq:boundary-link-reversal}
\end{equation}
where the operator definition and factorization theorem justify the relation.
The ultraviolet and rapidity kernels are properties of the operator
representation, not of the sign of the absorptive boundary.  Therefore a
common evolution map preserves \cref{eq:boundary-link-reversal}.  The process
map may still select different hard/color contractions.


\section{Two-scale TMD evolution as a geometric flow}
\label{sec:two-scale-evolution}

\subsection{Convention-consistent evolution equations}

For a renormalized TMD coefficient with complete operator identity
\(\mathfrak o\), define
\begin{align}
 \frac{\dd}{\dd\ln\mu}
 \ln \widetilde F_{\mathfrak o}(x,\bM;\mu,\zeta)
 &=
 \gamma_F^{a}(\mu,\zeta),
 \label{eq:mu-evolution}\\
 \frac{\dd}{\dd\ln\sqrt\zeta}
 \ln \widetilde F_{\mathfrak o}(x,\bM;\mu,\zeta)
 &=
 -\cD_a(\bM;\mu).
 \label{eq:zeta-evolution}
\end{align}
We choose the normalization
\begin{align}
 \frac{\partial\gamma_F^a}{\partial\ln\sqrt\zeta}
 &=-\Gamma_a(\alpha_s),
 \label{eq:gamma-zeta-consistency}\\
 \frac{\dd\cD_a}{\dd\ln\mu}
 &=\Gamma_a(\alpha_s),
 \label{eq:d-cusp-consistency}
\end{align}
so mixed derivatives commute exactly.  In this convention,
\begin{equation}
 \gamma_F^a(\mu,\zeta)
 =\gamma_a(\alpha_s(\mu))
 -\Gamma_a(\alpha_s(\mu))
  \ln\frac{\sqrt\zeta}{\mu}.
 \label{eq:gamma-f-form}
\end{equation}
Literature conventions that use \(\gamma_K=2\Gamma_a\), derivatives with
respect to \(\ln\mu^2\), or a Collins--Soper kernel \(\widetilde K=-\cD\) are
related by a stored convention adapter.  Factor-of-two differences may never
be resolved by inspecting numerical values alone.

\begin{theorem}[Exact path independence]
Assume \cref{eq:gamma-zeta-consistency,eq:d-cusp-consistency} and a simply
connected scale domain free of singularities.  The exact evolution factor
between \((\mu_i,\zeta_i)\) and \((\mu_f,\zeta_f)\) is independent of the path
in the \((\ln\mu,\ln\sqrt\zeta)\) plane.
\end{theorem}

\begin{proof}
Introduce the one-form
\begin{equation}
 \omega_a
 =\gamma_F^a(\mu,\zeta)\,\dd\ln\mu
 -\cD_a(\bM;\mu)\,\dd\ln\sqrt\zeta.
 \label{eq:evolution-one-form}
\end{equation}
Its exterior derivative is
\begin{equation}
 \dd\omega_a
 =-
 \left[
 \frac{\partial\gamma_F^a}{\partial\ln\sqrt\zeta}
 +\frac{\partial\cD_a}{\partial\ln\mu}
 \right]
 \dd\ln\mu\wedge\dd\ln\sqrt\zeta=0.
 \label{eq:evolution-closed-form}
\end{equation}
Thus \(\omega_a\) is closed and locally exact.  Its contour integral depends
only on the endpoints.
\end{proof}

The exact solution is
\begin{equation}
 \widetilde F_f
 =R_a[\cC;\bM]
 \widetilde F_i,
 \qquad
 \ln R_a[\cC;\bM]
 =\int_{\cC}\omega_a.
 \label{eq:contour-evolution-factor}
\end{equation}

\subsection{Finite-order curl defect}

At truncated perturbative order, independently truncated
\(\gamma_F^{[N]}\) and \(\cD^{[N]}\) need not satisfy exact integrability.
Define
\begin{equation}
 \Omega_a^{[N]}(\mu,\bM)
 =
 \frac{\partial\gamma_F^{a,[N]}}{\partial\ln\sqrt\zeta}
 +
 \frac{\partial\cD_a^{[N]}}{\partial\ln\mu}.
 \label{eq:finite-order-curl}
\end{equation}
For two contours \(\cC_1,\cC_2\) with the same endpoints,
\begin{equation}
 \ln\frac{R_a^{[N]}[\cC_1]}{R_a^{[N]}[\cC_2]}
 =\int_{\Sigma(\cC_1-\cC_2)}
 \Omega_a^{[N]}
 \dd\ln\mu\,\dd\ln\sqrt\zeta.
 \label{eq:contour-defect}
\end{equation}
The defect is a perturbative truncation diagnostic.  It must not be hidden by
choosing one contour and declaring path independence by construction.

\begin{requirement}[Evolution implementation record]
Every evolved result stores the contour or prescription, perturbative orders
of all anomalous dimensions and the beta function, endpoint scales, profile
parameters, treatment of the finite-order curl, and numerical integration
tolerances.
\end{requirement}

\subsection{Admissible evolution prescriptions}

The formalism permits several implementations:
\begin{enumerate}
\item \textbf{Direct contour evolution:} integrate \cref{eq:evolution-one-form}
      on a specified contour and report the curl defect.
\item \textbf{Integrability-improved evolution:} modify one truncated
      anomalous dimension by a formally higher-order term so the finite-order
      field is conservative.
\item \textbf{\(\zeta\)-prescription:} choose a null-evolution curve
      \(\zeta=\zeta_\mu(\bM)\) on which the boundary distribution is separated
      from double-scale evolution, then evolve from that curve to the final
      point \cite{ScimemiVladimirov2017,ScimemiVladimirov2018}.
\item \textbf{Equivalent SCET/RRG path:} use \((\mu,\nu)\) scales with an
      explicit finite conversion to the project \((\mu,\zeta)\) scheme
      \cite{Chiu2012}.
\end{enumerate}
The choice is a scheme/prescription decision shared by the full calculation,
not a per-TMD fit option.

\subsection{Target and polarization independence}

The anomalous dimensions in
\cref{eq:mu-evolution,eq:zeta-evolution} are properties of the renormalized
operator.  At leading power they do not depend on whether the external target
is a proton, neutron, deuteron, or another hadron, nor on its
\(U,L,T,LL,LT,TT\) polarization.  They also do not depend on the named TMD
projection when the underlying operator and color representation are the
same.  Polarization dependence enters through:
\begin{itemize}
\item the matched boundary matrix element;
\item the small-\(\bM\) coefficient and collinear operator channel;
\item the transverse-rank projector;
\item process-specific angular and hard kernels.
\end{itemize}

\begin{requirement}[No polarization-specific CS kernel by default]
A separate Collins--Soper kernel for every spin-1 TMD is forbidden unless a
new operator definition or a controlled factorization-breaking mechanism is
identified.  Data tension must first be tested against boundary, matching,
power-correction, and process kernels.
\end{requirement}

\section{The Collins--Soper kernel across perturbative and nonperturbative domains}
\label{sec:cs-kernel}

\subsection{Small-\texorpdfstring{\(\bM\)}{bTMD} expansion}

For \(b=\abs{\bM}\ll\Lambda_{\QCD}^{-1}\),
\begin{equation}
 \cD_a(b;\mu)
 =\cD_a^{\pert}
 \left(L_b,\alpha_s(\mu)\right),
 \qquad
 L_b=\ln\frac{\mu^2b^2}{4e^{-2\gamma_E}},
 \label{eq:perturbative-d-kernel}
\end{equation}
with coefficients constrained by
\cref{eq:d-cusp-consistency}.  The perturbative series is known to high order,
and the rapidity anomalous dimension has been calculated through four loops
\cite{MoultZhuZhu2022,DuhrMistlbergerVita2022}.  The implementation records the
actual coefficient order rather than inferring it from a resummation label.

\subsection{Large-\texorpdfstring{\(b\)}{b} completion}

The kernel is generically nonperturbative at large \(b\).  A general matched
representation is
\begin{equation}
 \cD_a(b;\mu)
 =
 \cD_a^{\pert}(b_{\rm pert}(b);\mu)
 +g_{K,a}(b)
 +\delta\cD_{a}^{\rm match}(b),
 \label{eq:large-b-kernel-decomposition}
\end{equation}
where \(b_{\rm pert}\) is a declared profile or regulated coordinate,
\(g_{K,a}\) is the nonperturbative remainder, and
\(\delta\cD^{\rm match}\) enforces smoothness or step-scaling consistency.
A \(b_*\) map is one admissible realization, not part of the definition of a
TMD.

The profile must satisfy at least
\begin{equation}
 b_{\rm pert}(b)=b+\cO(b^3/b_{\max}^2)
 \quad(b\to0),
 \qquad
 b_{\rm pert}(b)\le b_{\max},
 \label{eq:b-profile-limits}
\end{equation}
and its induced power corrections are a named uncertainty axis.

\subsection{Color-representation dependence}

Quark and gluon kernels are separate objects,
\begin{equation}
 \cD_q\ne\cD_g
 \label{eq:quark-gluon-d-separation}
\end{equation}
in general.  Quadratic-Casimir scaling is useful through the orders where it
is verified, while four-loop results involve generalized Casimir structures
\cite{MoultZhuZhu2022,DuhrMistlbergerVita2022}.  The model therefore stores
independent quark and gluon perturbative coefficients and does not impose a
nonperturbative relation
\(g_{K,g}=(C_A/C_F)g_{K,q}\) without evidence.

For independent gluon \(f\)- and \(d\)-type TMD operator channels, the
standard leading-power rapidity kernel is associated with the adjoint Wilson
representation.  A color-channel difference in the boundary or process hard
map does not by itself imply a different universal CS kernel.  If a process
requires a matrix-valued soft anomalous dimension, it is no longer represented
as two independent universal gluon TMDs without an explicit generalized
factorization theorem.

\subsection{Calibration information}

The nonperturbative kernel may be constrained by:
\begin{itemize}
\item global Drell--Yan, vector-boson, and SIDIS analyses;
\item lattice quasi-TMD or TMD-wave-function determinations after controlled
      matching;
\item the microscopic model only if it includes the same soft operator and
      subtracts overlap with the boundary phase budget;
\item common fits using multiple target and polarization channels.
\end{itemize}
Recent lattice calculations provide direct constraints on the quark
Collins--Soper kernel over nonperturbative transverse separations
\cite{Avkhadiev2024,Bollweg2024,TanEtAl2025}.  They are calibration or
validation data in a declared scheme; they are not silently copied into the
gluon kernel.

\begin{requirement}[Kernel universality test]
A common quark or gluon kernel must describe all supported target and
polarization channels after their distinct boundaries and coefficient
functions are included.  Failure is reported before introducing a
polarization-dependent kernel.
\end{requirement}

\section{Rank-aware Fourier--Bessel representation}
\label{sec:rank-transforms}

\subsection{Irreducible transverse harmonics}

A coefficient with \(SO(2)\) weight \(m\) multiplies an irreducible tensor
\(\mathcal T_m(\hat\kT)\).  With the project Fourier convention,
\begin{equation}
 \widetilde F^{(m)}(x,b)
 =2\pi\,i^{m}\int_0^\infty
 k\,\dd k\,
 J_{|m|}(bk)\,
 N_m^{k}(k;M_{\rm ref})
 F^{(m)}(x,k),
 \label{eq:rank-forward-bessel}
\end{equation}
where \(N_m^k\) includes the extracted powers and sign convention stored in
\(\mathfrak r\).  The inverse is
\begin{equation}
 F^{(m)}(x,k)
 =\frac{i^{-m}}{2\pi}
 \int_0^\infty b\,\dd b\,
 J_{|m|}(bk)\,
 N_m^b(b;M_{\rm ref})
 \widetilde F^{(m)}(x,b).
 \label{eq:rank-inverse-bessel}
\end{equation}
The normalization factors obey the registry's round-trip identity, not a
universal scalar formula.

\subsection{Evolution does not change transverse rank}

For a scalar evolution factor \(R_a(b)\) associated with one operator
representation,
\begin{equation}
 \widetilde F_f^{(m)}(x,b)
 =R_a(b)\widetilde F_i^{(m)}(x,b).
 \label{eq:rank-preserving-evolution}
\end{equation}
It does not mix \(m\) or target irreducible channels.  Mixing can occur in the
small-\(b\) collinear operator basis or process contraction, not through the
universal radial evolution factor.

\begin{proposition}[Rank/evolution commutation]
If the evolution kernel is rotationally invariant in the transverse plane and
multiplicative in \(b\) space, then evolution commutes with the projection onto
each irreducible \(SO(2)\) harmonic.
\end{proposition}

\begin{proof}
The kernel depends only on \(b=\abs{\bM}\) and therefore commutes with the
rotation representation and its irreducible projectors.  Multiplication by
\(R_a(b)\) preserves the harmonic weight.  The subsequent Bessel transform
uses the same \(J_{|m|}\) kernel.
\end{proof}

\subsection{Derivative moments}

Weighted transverse moments can be represented by derivatives at small
\(b\), but these derivatives are ultraviolet sensitive and require operator
matching.  Schematically,
\begin{equation}
 F^{(n)}(x)
 \sim
 \left.
 \left(-\frac{\partial}{\partial b^2}\right)^n
 \widetilde F(x,b)
 \right|_{b\to0}^{\ren}.
 \label{eq:derivative-moment}
\end{equation}
The superscript \(\ren\) is essential.  A finite-grid numerical derivative of
a model curve is not automatically a renormalized twist-3 or higher-twist
matrix element.

\subsection{Numerical transform contract}

Every transform manifest records:
\begin{itemize}
\item harmonic and rank;
\item forward and inverse phases;
\item extracted mass powers;
\item integration range and quadrature rule;
\item small- and large-argument Bessel treatment;
\item interpolation basis;
\item round-trip residuals on analytic and production test functions;
\item tail estimates beyond the finite grids.
\end{itemize}
A transform cannot pass merely because the same implementation is used in
both directions; at least one independent analytic oracle is required.

\section{Small-\texorpdfstring{\(b_{\mathrm{TMD}}\)}{bTMD} operator product expansion}
\label{sec:small-b-ope}

\subsection{General typed expansion}

For a renormalized TMD coefficient \(F_A^{a/h}\),
\begin{equation}
 \widetilde F_A^{a/h}(x,b;\mu,\zeta)
 =
 \sum_{B,j}
 \left[
 C_{A\leftarrow B}^{a\leftarrow j}
 (b;\mu,\zeta;\mathfrak s)
 \otimesx
 f_B^{j/h}(\mu)
 \right](x)
 +\cO\!\left((b\Lambda_{\QCD})^{p_A}\right).
 \label{eq:generic-small-b-ope}
\end{equation}
Here \(B\) labels a complete collinear operator basis, including twist,
polarization, flavor, charge-conjugation, total-derivative, and multiparton
information.  The coefficient is target independent; the proton, neutron,
deuteron, and tensor-polarized information resides in the collinear matrix
element.

For nonzero rank,
\begin{equation}
 \widetilde\Phi_A^{i_1\cdots i_r}
 =
 b^{\langle i_1}\cdots b^{i_r\rangle}
 \sum_{B,j}
 C_{A\leftarrow B}^{(r),a\leftarrow j}
 \otimesx f_B^{j/h}
 +\cdots,
 \label{eq:ranked-small-b-ope}
\end{equation}
with the extracted powers and reference mass governed by \(\mathfrak r\).

\subsection{Twist-2 channels}

The leading-twist small-\(b\) basis includes, where allowed by quantum
numbers:
\begin{itemize}
\item unpolarized quark and gluon PDFs;
\item helicity quark and gluon PDFs;
\item quark transversity;
\item tensor-polarized rank-zero quark and gluon operators for a spin-1
      target;
\item linearly polarized gluon TMD matching onto ordinary collinear partons;
\item special spin-1 double-helicity-flip gluon operators where a collinear
      twist-2 matrix element exists.
\end{itemize}
The generic twist-2 matching structure and its allowed polarization channels
have been worked out at fixed order \cite{GutierrezReyes2017}; linearly
polarized gluon matching is known to high perturbative order
\cite{GutierrezReyes2019,LuoYangZhuZhu2019,ZhuLinearlyPolarized2025}, and
transversity matching has likewise advanced to high order
\cite{ZhuTransversity2025,ZhuDGLAP2026}.

At tree level for the rank-zero unpolarized-like spin-1 channel,
\begin{equation}
 \widetilde f_{1LL}^{a/D}(x,b)
 =f_{1LL}^{a/D}(x)+\cO(\alpha_s,b^2\Lambda_{\QCD}^2),
 \label{eq:f1ll-tree-matching}
\end{equation}
provided the same \(LL\) convention is used on both sides.  Equivalently, the
helicity-state tensor difference \(\delta_T f\) can be used as the
convention-stable collinear coordinate.

The fact that a fixed-order coefficient vanishes does not imply that the
microscopic boundary function is exactly zero.  It can indicate that the
matching begins at a higher order or that the leading collinear operator has a
different twist.

\subsection{Twist-3 and naive-\texorpdfstring{\(T\)-odd}{T-odd} channels}

Sivers, Boer--Mulders, and several worm-gear structures have small-\(b\)
expansions involving quark--gluon--quark or tri-gluon collinear operators
\cite{ScimemiVladimirovTwist3}.  In a schematic project convention,
\begin{align}
 \widetilde f_{1T}^{\perp q[\eta]}
 &\sim
 \eta\,
 C_{\rm Siv}^{q\leftarrow j}
 \otimesx T_F^j
 +C_{\rm Siv}^{q\leftarrow G}
 \otimesx T_G
 +\cdots,
 \label{eq:sivers-twist3-matching}\\
 \widetilde h_1^{\perp q[\eta]}
 &\sim
 \eta\,
 C_{\rm BM}^{q\leftarrow j}
 \otimesx T_F^{(\sigma),j}
 +\cdots.
 \label{eq:boer-mulders-twist3-matching}
\end{align}
The displayed sign is schematic: the exact sign, factors of \(M\), Fourier
phase, and definition of \(T_F\) are stored in the coefficient registry.  A
future/past sign may never be imposed twice---once in the microscopic Wilson
boundary and again in the matching coefficient.

Gluon naive-\(T\)-odd functions match onto independent \(f\)-type and
\(d\)-type tri-gluon operators, with quark--gluon mixing where allowed:
\begin{equation}
 \begin{pmatrix}
  \widetilde F_g^{(f)}\\
  \widetilde F_g^{(d)}
 \end{pmatrix}
 =
 \begin{pmatrix}
  C_{ff}&C_{fd}&C_{fq}\\
  C_{df}&C_{dd}&C_{dq}
 \end{pmatrix}
 \otimesx
 \begin{pmatrix}
  T_G^{(f)}\\
  T_G^{(d)}\\
  T_F^{q}
 \end{pmatrix}
 +\cdots.
 \label{eq:gluon-todd-matching-matrix}
\end{equation}
No universal scalar ``gluon Sivers coefficient'' replaces this matrix.

\subsection{Spin-1 tensor channels}

Target tensor polarization does not alter the short-distance coefficient of a
given QCD operator, but it enlarges the set of nonzero external matrix
elements.  Each spin-1 TMD is classified by:
\begin{equation}
 \left(
 \text{parton operator},\,
 \text{target irrep},\,
 r,\,
 \text{link parity},\,
 \text{collinear twist},\,
 \text{lowest nonzero order}
 \right).
 \label{eq:spin1-ope-classification}
\end{equation}
A rank-zero \(LL\) distribution can share the unpolarized twist-2 coefficient
while using a tensor-polarized collinear matrix element.  An \(LT\) or \(TT\)
coefficient may instead have zero unweighted integral and match to a weighted
moment, a helicity-flip operator, or a higher-twist multiparton operator.  The
registry must say which case applies; it cannot infer the collinear anchor
from the target-channel name.

\subsection{Power corrections and overlap domain}

The OPE is accepted only in a domain where
\begin{equation}
 b\Lambda_{\QCD}\ll1,
 \qquad
 b\Lambda_{\LF}\ll1
 \quad\text{or with a supplied LF-to-QCD step-scaling map}.
 \label{eq:ope-domain}
\end{equation}
The residual is decomposed as
\begin{equation}
 \delta_{\rm OPE}
 =\delta_{b^2\Lambda^2}
 +\delta_{\rm twist}
 +\delta_{\alpha_s}
 +\delta_{\rm basis}
 +\delta_{\rm LF\to QCD}.
 \label{eq:ope-error-budget}
\end{equation}
Changing the large-\(b\) model cannot be used to hide failure of small-\(b\)
matching.


\section{Collinear evolution and operator mixing}
\label{sec:collinear-evolution}

\subsection{Nonsinglet evolution}

For a nonsinglet operator coordinate \(f_{\NS}\),
\begin{equation}
 \frac{\dd}{\dd\ln\mu^2}f_{\NS}(x;\mu)
 =
 \left[P_{\NS}(\alpha_s)\otimesx f_{\NS}\right](x;\mu).
 \label{eq:nonsinglet-dglap}
\end{equation}
Different operator polarizations use different splitting functions: the
unpolarized, helicity, and transversity kernels are distinct.  The target
polarization does not change the kernel of a fixed operator.  Thus a tensor
helicity-state difference of the unpolarized quark operator evolves with the
same unpolarized nonsinglet kernel as the corresponding target average.

\subsection{Singlet quark--gluon system}

Define the flavor singlet
\begin{equation}
 \Sigma(x;\mu)=\sum_{q=1}^{n_f}
 \left[q(x;\mu)+\bar q(x;\mu)\right].
 \label{eq:singlet-definition}
\end{equation}
For each operator polarization that permits quark--gluon mixing,
\begin{equation}
 \frac{\dd}{\dd\ln\mu^2}
 \begin{pmatrix}
  \Sigma\\g
 \end{pmatrix}
 =
 \begin{pmatrix}
  P_{qq}&P_{qg}\\
  P_{gq}&P_{gg}
 \end{pmatrix}
 \otimesx
 \begin{pmatrix}
  \Sigma\\g
 \end{pmatrix}.
 \label{eq:singlet-evolution}
\end{equation}
For a spin-1 target, the same matrix acts on the tensor-polarized singlet
coordinates
\begin{equation}
 \delta_T\bm f
 =
 \begin{pmatrix}
  \delta_T\Sigma\\
  \delta_T g
 \end{pmatrix},
 \label{eq:tensor-singlet-vector}
\end{equation}
because the local twist-2 operators are the same and only the external target
matrix element changes.  This statement does not extend automatically to
ranked or helicity-flip tensor TMDs whose collinear operator basis differs.

\begin{proposition}[Momentum-sum preservation]
If the splitting matrix satisfies the standard first-moment momentum
constraints and the threshold maps are consistent, then the total plus
momentum
\(
\int_0^1\dd x\,x[\Sigma(x)+g(x)]
\)
is invariant under \cref{eq:singlet-evolution}.
\end{proposition}

\begin{proof}
Take the first \(x\)-moment of \cref{eq:singlet-evolution}.  Associativity of
Mellin convolution converts it to multiplication by the first moments of the
splitting functions.  The left null vector expressing energy--momentum
conservation annihilates the splitting matrix.
\end{proof}

\subsection{Helicity, transversity, and double-helicity-flip sectors}

Helicity evolution uses the polarized splitting matrix
\(\Delta P_{ij}\).  Quark transversity is chiral odd and has no ordinary
twist-2 gluon mixing for a spin-\(1/2\) target.  A spin-1 target can support
independent gluon double-helicity-flip operators, but they are not inserted
into the quark transversity equation by analogy.  Their evolution is defined
in the operator basis appropriate to their helicity transfer and color
representation.

The implementation stores separate evolution families:
\begin{equation}
 \mathfrak e\in
 \left\{
 \mathrm{unpolarized},\,
 \mathrm{helicity},\,
 \mathrm{transversity},\,
 \mathrm{gluon\ double\ flip},\,
 \mathrm{twist\!-\!3\ multiparton}
 \right\}.
 \label{eq:collinear-evolution-families}
\end{equation}
A common interface does not mean a common splitting kernel.

\subsection{Twist-3 multiparton evolution}

The twist-3 functions entering
\cref{eq:sivers-twist3-matching,eq:boer-mulders-twist3-matching,eq:gluon-todd-matching-matrix}
are functions of more than one momentum fraction before special diagonal
projections are taken.  Their evolution is an integral operator on a vector
space of quark--gluon--quark and tri-gluon correlators,
\begin{equation}
 \frac{\dd}{\dd\ln\mu^2}\bm T(\bm x;\mu)
 =\bm K_{\tw3}(\mu)\otimes_{\bm x}\bm T(\bm x;\mu).
 \label{eq:twist3-evolution}
\end{equation}
Evolving only the diagonal Qiu--Sterman function while ignoring generated
off-diagonal support is an approximation that requires its own closure test.
The twist-3 operator vector retains \(f/d\) color, flavor, charge-conjugation,
and link-parity labels.

\subsection{Heavy-flavor thresholds}

At a threshold \(\mu_h\), the \(n_f\)- and \(n_f+1\)-flavor theories are
related by
\begin{equation}
 \bm f^{(n_f+1)}(\mu_h)
 =\bm A_{n_f+1\leftarrow n_f}(\mu_h)
 \otimesx
 \bm f^{(n_f)}(\mu_h),
 \label{eq:collinear-threshold-map}
\end{equation}
and the TMD matching coefficients, anomalous dimensions, hard functions, and
running coupling must change consistently.  Individual scheme-dependent
parton functions need not be exactly continuous, but physical observables and
conserved moments must be continuous to the retained order.

A threshold record contains:
\begin{equation}
 \mathfrak h=
 \left(
 m_h,\,\mu_h,\,n_f^{\rm in},\,n_f^{\rm out},\,
 \bm A,\,\alpha_s\text{ matching},\,
 \TMD\text{ coefficient matching},\,
 \text{order}
 \right).
 \label{eq:threshold-record}
\end{equation}
The flavor count cannot change merely because the final hard scale crosses a
mass value.

\subsection{Collinear/TMD evolution consistency}

The small-\(b\) OPE and evolution must satisfy
\begin{equation}
 \cE_{\TMD}
 \left[
 \bm C(\mu_i,\zeta_i)\otimesx\bm f(\mu_i)
 \right]
 =
 \bm C(\mu_f,\zeta_f)
 \otimesx
 \cE_{\rm coll}\left[\bm f(\mu_i)\right]
 +\cO(\delta_{\rm pert}),
 \label{eq:ope-evolution-commutation}
\end{equation}
where all coefficient, anomalous-dimension, and beta-function orders are
compatible.  This is a stronger test than evolving the TMD boundary and PDF
independently and comparing only at one scale.

\section{Commutation with the microscopic and nuclear hierarchy}
\label{sec:nuclear-evolution}

\subsection{Evolution acts on operators, not on target labels}

Let \(\cN\) be a linear nuclear composition map from nucleon operator matrix
elements to a deuteron matrix element.  For a one-body impulse term,
\begin{equation}
 \left[\cN F\right](x)
 =\int_x^1\frac{\dd y}{y}\,
 S(y)F\left(\frac{x}{y}\right),
 \label{eq:nuclear-mellin-convolution}
\end{equation}
with spin and transverse variables suppressed.  Let \(\cE\) be a linear
collinear or TMD evolution operator whose kernel is target independent.

\begin{theorem}[Evolution--impulse commutation]
If the nuclear spectral kernel \(S\) is independent of \(\mu\) and \(\zeta\)
in the selected operator representation, the integrals converge, and the
same partonic operator basis is used, then
\begin{equation}
 \cE\circ\cN=\cN\circ\cE.
 \label{eq:evolution-nuclear-commutation}
\end{equation}
\end{theorem}

\begin{proof}
Both \(\cE\) and \(\cN\) are Mellin convolutions in the longitudinal
fraction.  Associativity and Fubini's theorem allow the integration order to
be exchanged.  TMD rapidity evolution multiplies the nucleon TMD by a radial
factor in \(b\) space and therefore also commutes with the nuclear
longitudinal convolution and the nuclear transverse phase.  The conclusion
fails if the nuclear kernel itself contains unresolved partonic scale
dependence or if different operator bases are used.
\end{proof}

For the spin-resolved impulse parent,
\begin{equation}
 \widetilde F_{a/D,I}^{\IA}(x,b;\mu,\zeta)
 =
 \sum_N\int\frac{\dd y}{y}\dd^2\pT\,
 S_{I}^{N/D}(y,\pT)
 e^{i(x/y)\bM\cdot\pT}
 \widetilde F_{a/N}
 \left(\frac{x}{y},b;\mu,\zeta\right),
 \label{eq:evolved-impulse-convolution}
\end{equation}
the universal evolution factor can be applied before or after the integration
when the assumptions of the theorem hold.

\subsection{When commutation fails}

The theorem does not license one universal nucleon evolution factor for every
nuclear mechanism.  Separate matched operator bases are required for:
\begin{itemize}
\item explicit pion and \(\Delta\) partonic operators;
\item coherent diffractive multi-nucleon operators;
\item short-distance six-quark or hidden-color operators;
\item exchange-current operators with their own anomalous dimensions;
\item factorization regimes in which nuclear Glauber exchange couples the two
      hadrons or measured final state.
\end{itemize}
In those cases the full operator vector evolves,
\begin{equation}
 \frac{\dd}{\dd\ln\mu}
 \begin{pmatrix}
 \cO_{(1)}\\\cO_{(2)}\\\cO_{\pi}\\\cO_{\Delta}\\\cO_{\shortd}
 \end{pmatrix}
 =\bm\gamma_{\nuc}
 \begin{pmatrix}
 \cO_{(1)}\\\cO_{(2)}\\\cO_{\pi}\\\cO_{\Delta}\\\cO_{\shortd}
 \end{pmatrix},
 \label{eq:nuclear-operator-mixing}
\end{equation}
with subtraction terms preventing the same short-distance physics from being
carried by multiple sectors.

\subsection{Microscopic member identity}

A microscopic member \(m\) must remain the same through matching, evolution,
nuclear composition, and process assembly:
\begin{equation}
 \theta_m
 \longrightarrow
 \bm W_m^{\LF}
 \longrightarrow
 \bm W_m^{\QCD}
 \longrightarrow
 \bm F_m(\mu,\zeta)
 \longrightarrow
 \sigma_m.
 \label{eq:member-identity-chain}
\end{equation}
One may not combine the quark boundary from member \(m\), gluon boundary from
member \(m'\), and nuclear state from member \(m''\) unless the ensemble
construction explicitly declares them independent nuisance axes rather than
one correlated microscopic state.

\begin{requirement}[No post-evolution reassembly]
Evolution acts on the resolved component graph.  Proton, neutron, flavor,
parton species, target channel, rank, link/color class, nuclear mechanism,
and member labels remain available after evolution.  Summing them before
evolution is allowed only when the evolution kernels and matching maps are
proven identical on the summed subspace.
\end{requirement}

\section{Boundary initialization and the separation of physics}
\label{sec:boundary}

\subsection{Reference boundary}

Choose a reference point or curve
\begin{equation}
 \mathfrak b_0=
 \left(\mu_0,\zeta_0,\mathfrak s_{\TMD},r_{\LF}\right)
 \label{eq:boundary-record}
\end{equation}
in a domain where the Hamiltonian matrix element, regulator-to-QCD matching,
and perturbative evolution are all controlled.  The initial TMD is
\begin{equation}
 \widetilde{\bm F}_0
 =
 \bm Z_{\LF\to\QCD}(\mathfrak b_0)
 \otimesx\widetilde{\bm W}_{\LF}
 +\widetilde{\bm\delta}_{\rm miss}
 +\widetilde{\bm\delta}_{\rm trunc}.
 \label{eq:initialized-boundary}
\end{equation}
The reference scale is not required to equal the hadron model cutoff, the
small-\(b\) matching scale, or the final hard scale.

\subsection{No independent Gaussian per function}

A numerical regulator or basis may produce approximately Gaussian transverse
profiles.  This does not justify replacing each microscopic projection by
\begin{equation}
 F_A(x,k_T)\propto f_A(x)
 \exp[-k_T^2/\lambda_A^2]
 \label{eq:forbidden-independent-gaussian}
\end{equation}
with one freely fitted \(\lambda_A\) per named TMD.  Transverse structure must
come from the common state, matching operator, or a small shared discrepancy
basis whose parameters affect all related projections.

Admissible discrepancy operators include:
\begin{itemize}
\item a common unresolved high-Fock transverse form factor;
\item a representation-specific large-\(b\) kernel;
\item a controlled missing-current operator;
\item a shared quark--gluon vertex correction;
\item a common nuclear short-distance operator.
\end{itemize}
Each has an operator identity, power counting, and withheld test.

\subsection{Phase budget from Volume III}

The link-odd boundary inherits the phase budget
\begin{equation}
 \operatorname{Im}\ln W_{\ren}
 =
 \operatorname{Im}\ln W_{\unsub}
 -\frac12\operatorname{Im}\ln S
 +\operatorname{Im}\ln Z_{\UV}
 +\operatorname{Im}\ln R_{\rap}.
 \label{eq:phase-budget-repeated}
\end{equation}
A boundary-only rescattering calculation and a joint soft-sector calculation
are mutually exclusive representations of the same soft exchange unless an
explicit overlap subtraction is supplied.  The CS kernel cannot be fitted to
reproduce a phase already included in the microscopic boundary.

\subsection{Positivity after matching and evolution}

At the regulated forward amplitude level, Gram positivity can be structural.
A finite-order soft-subtracted TMD in a particular scheme is not guaranteed to
remain pointwise positive.  Therefore:
\begin{itemize}
\item positivity is tested at every layer where a theorem applies;
\item negative values generated by perturbative matching are not clipped;
\item physical cross sections and density-matrix combinations are tested in
      the domain of factorization;
\item a positivity failure is classified as state, matching, truncation,
      scheme, or numerical, not repaired silently.
\end{itemize}

\section{Process-factorization records}
\label{sec:process-records}

\subsection{Typed process object}

A process factorization record is
\begin{equation}
 \begin{aligned}
 \mathfrak P=\bigl(&
 \text{external states},\,
 \text{hard channel},\,
 \text{measurement},\,
 \text{TMD/FF basis},\,
 \text{link map},\,
 \text{color map},\\
 &\text{soft prescription},\,
 \text{power counting},\,
 \text{Glauber status},\,
 \text{factorization status}
 \bigr).
 \end{aligned}
 \label{eq:process-record}
\end{equation}
The factorization status is one of
\begin{equation}
 \texttt{ESTABLISHED},\quad
 \texttt{CONDITIONAL},\quad
 \texttt{EXPLORATORY},\quad
 \texttt{BROKEN},\quad
 \texttt{UNASSESSED}.
 \label{eq:factorization-statuses}
\end{equation}
A \texttt{BROKEN} or \texttt{UNASSESSED} object cannot consume universal
separate-hadron TMDs as an authoritative prediction.

\subsection{Drell--Yan}

For a color-singlet Drell--Yan channel at \(q_T\ll Q\), a rank-resolved
structure function is schematically
\begin{align}
 W_m^{\DY}(q_T,Q)
 ={}&
 \sum_{a,b}
 H_{ab,m}^{\DY}(Q,\mu)
 \int_0^\infty\frac{b\,\dd b}{2\pi}
 J_m(bq_T)
 \nonumber\\
 &\times
 \widetilde F_{a/h_1}^{[\gamma_-],(r_1)}
 (x_1,b;\mu,\zeta_1)
 \widetilde F_{b/h_2}^{[\gamma_-],(r_2)}
 (x_2,b;\mu,\zeta_2)
 \cM_m(b),
 \label{eq:dy-w-term}
\end{align}
where \(\zeta_1\zeta_2\) is tied to \(Q^4\) in the chosen convention and
\(\cM_m\) contains the declared mass and angular factors.  The deuteron target
channel is a projector on its spin-1 density matrix, not a change of the
universal evolution kernel.

\subsection{SIDIS}

For SIDIS,
\begin{align}
 W_m^{\SIDIS}(P_{hT},Q)
 ={}&
 \sum_q e_q^2 H_{q,m}^{\SIDIS}(Q,\mu)
 \int_0^\infty\frac{b\,\dd b}{2\pi}
 J_m\!\left(\frac{bP_{hT}}{z_h}\right)
 \nonumber\\
 &\times
 \widetilde F_{q/D}^{[\gamma_+],(r_F)}
 (x_B,b;\mu,\zeta_F)
 \widetilde D_{h/q}^{(r_D)}
 (z_h,b;\mu,\zeta_D)
 \cM_m^{\SIDIS}(b,z_h).
 \label{eq:sidis-w-term}
\end{align}
The fragmentation-function scheme, rank convention, and Fourier scaling by
\(z_h\) are part of the process record.  The future-pointing link is not
selected from the sign of a fitted Sivers curve; it is supplied by the
factorization theorem.

\subsection{Inclusive DIS}

Inclusive DIS uses collinear factorization rather than a TMD product.  The
rank-zero tensor relation
\begin{equation}
 b_1(x,Q^2)
 =\frac12\sum_q e_q^2
 \left[
 \delta_T q(x,Q^2)+\delta_T\bar q(x,Q^2)
 \right]
 +\cO(\alpha_s,Q^{-2})
 \label{eq:b1-evolved-relation}
\end{equation}
provides a direct check that the TMD small-\(b\) matching, collinear evolution,
and tensor sign adapter are mutually consistent.  Integrating a SIDIS TMD
formula is not used as a substitute for the inclusive hard coefficients.

\subsection{Gluon-sensitive channels}

A gluon process can require a vector of independent color/link channels,
\begin{equation}
 W_m^{g,\proc}
 =
 \sum_{c=f,d}
 H_{m,c}^{g,\proc}
 \otimes
 \widetilde F_{g/D}^{(c)}
 \otimes
 \widetilde X_{\proc}^{(c)}.
 \label{eq:gluon-color-process-map}
\end{equation}
The hard factors determine the \(f/d\) weights.  For colored final states,
Glauber exchange or color entanglement can invalidate the product form.  The
process record must either cite the required factorization result, carry a
conditional generalized soft matrix, or remain exploratory.

\begin{requirement}[Process qualification before prediction]
A named process must supply its operator basis, link directions, color
weights, hard factors, measurement kernel, scale relation, fixed-order
reference, and factorization/Glauber status before the model can produce an
authoritative observable.
\end{requirement}


\section{Matched low-, intermediate-, and high-transverse-momentum observables}
\label{sec:wy}

\subsection{Additive \texorpdfstring{\(W+Y\)}{W+Y} construction}

For a process and angular structure \(m\), define
\begin{equation}
 \frac{\dd\sigma_m}{\dd\cP}
 =W_m^{[N]}+Y_m^{[N]}+\delta_{m,\rm pow},
 \label{eq:wy-master}
\end{equation}
where \(W_m^{[N]}\) is the TMD-resummed expression at declared accuracy and
\begin{equation}
 Y_m^{[N]}
 =
 \sigma_{m,\FO}^{[N]}
 -\left[W_m^{[N]}\right]_{\asy,\FO}^{[N]}.
 \label{eq:y-definition}
\end{equation}
The asymptotic operator expands the same hard factors, matching coefficients,
evolution exponent, measurement variables, masses, schemes, and angular
projectors used in \(W_m\).  A fixed-order cross section and an asymptotic
term computed in different conventions cannot define a valid \(Y\) term.

\begin{proposition}[Fixed-order recovery]
If \(\left[W_m\right]_{\asy,\FO}^{[N]}\) is the complete singular expansion of
\(W_m\) through the same order as \(\sigma_{m,\FO}^{[N]}\), then
\(W_m+Y_m\) reproduces the fixed-order result through order \(N\) when
expanded in \(\alpha_s\), while retaining all-order logarithmic resummation at
small \(q_T\).
\end{proposition}

\begin{proof}
Expand \(W_m\) through order \(N\).  Its singular fixed-order part cancels
against the subtraction in \cref{eq:y-definition}, leaving
\(\sigma_{m,\FO}^{[N]}\).  Terms beyond order \(N\) are the resummed
contribution, and power corrections remain explicit.
\end{proof}

\subsection{Behavior in the two limits}

At \(q_T/Q\to0\), the \(Y\) term is nonsingular and should be parametrically
suppressed relative to the leading logarithmic content of \(W\).  At
\(q_T\sim Q\), the resummed term alone is not a controlled full description;
the sum must approach the fixed-order prediction within the declared
power-correction and matching uncertainty.

The model boundary is never retuned to absorb a deficient \(Y\) term.  A
large discrepancy at moderate \(q_T\) is diagnosed among:
\begin{itemize}
\item missing fixed-order channels;
\item inconsistent asymptotic subtraction;
\item resummation/profile choice;
\item kinematic power corrections;
\item factorization breakdown;
\item numerical transform or interpolation error.
\end{itemize}

\subsection{Profiled matching}

A smooth profile can be introduced to suppress resummation outside its natural
domain,
\begin{equation}
 \sigma_m^{\rm matched}
 =w(q_T/Q)W_m
 +\sigma_{m,\FO}
 -w(q_T/Q)\left[W_m\right]_{\asy,\FO},
 \label{eq:profiled-wy}
\end{equation}
with \(w(0)=1\) and \(w(q_T/Q\gtrsim1)\to0\).  The profile is part of the
accuracy record.  Variations of its transition location and smoothness probe
matching uncertainty, but a profile may not remove a physical fixed-order
channel or alter the expansion through the claimed order.

\subsection{Rank-resolved subtraction}

The asymptotic subtraction is performed separately for every harmonic and
tensor rank:
\begin{equation}
 Y_m
 =\sigma_{\FO,m}-[W_m]_{\asy,\FO}.
 \label{eq:rank-resolved-y}
\end{equation}
Using a rank-zero asymptotic term for a rank-one or rank-two modulation can
produce the correct dimensions after ad hoc mass factors while still being
physically wrong.  The `RankSpec' of the TMD, hard tensor, and fixed-order
projector must agree.

\subsection{Multiplicative alternatives}

Multiplicative or inverse-error-weighted matching can be studied as an
alternative,
\begin{equation}
 \sigma^{\rm mult}
 =W\,\frac{\sigma_{\FO}}{[W]_{\asy,\FO}},
 \label{eq:multiplicative-matching}
\end{equation}
but only in a domain where the denominator is nonzero and its perturbative
expansion is controlled.  It is not the project default.  Additive and
multiplicative results are alternative matching schemes, not uncertainty
members to be added simultaneously.

\section{Perturbative accuracy and resummation bookkeeping}
\label{sec:accuracy}

\subsection{The accuracy tuple is authoritative}

A label such as \(\mathrm{N}^3\mathrm{LL}'\) is insufficient because counting
conventions differ.  The authoritative record is
\begin{equation}
 \mathfrak A=
 \left(
 n_{\Gamma},n_{\gamma},n_{\cD},n_{\beta},
 n_C,n_H,n_P,n_{\FO},n_Y,
 \mathfrak p_{\rm evol},\mathfrak p_{\rm match}
 \right),
 \label{eq:accuracy-tuple}
\end{equation}
where the entries specify the perturbative order of the cusp, noncusp, and
rapidity anomalous dimensions, beta function, small-\(b\) coefficients, hard
factor, collinear splitting functions, fixed-order cross section, and
asymptotic subtraction, plus the evolution and matching prescriptions.

The human-readable resummation label is generated from \(\mathfrak A\) under
a versioned project convention.  A primed label means that specified boundary
or hard/matching coefficients are included one fixed order beyond the minimal
unprimed logarithmic counting; the exact entries are still listed.

\subsection{Scale families}

The calculation can contain:
\begin{equation}
 \mu_H,\quad
 \mu_B,\quad
 \mu_S,\quad
 \mu_b,\quad
 \mu_{\FO},\quad
 \zeta_1,\zeta_2,\quad
 \mu_{\rm thr},\quad
 \mu_{\LF\to\QCD}.
 \label{eq:scale-family}
\end{equation}
Not all are independent in every scheme.  A variation plan declares which
scales are varied, their correlations, constraints such as
\(\zeta_1\zeta_2=Q^4\), and regions excluded to avoid artificial large
logarithms.

Scale variation is not the only perturbative uncertainty.  The error budget
also includes:
\begin{itemize}
\item finite-order evolution curl;
\item OPE truncation and unknown coefficient orders;
\item heavy-threshold matching;
\item \(W+Y\) transition/profile choice;
\item fixed-order numerical uncertainty;
\item scheme-conversion truncation.
\end{itemize}

\subsection{No accuracy laundering}

A high-order cusp anomalous dimension does not upgrade a calculation whose
hard factor, matching coefficient, polarized splitting function, or \(Y\) term
is known only at lower order.  Every observable reports both the global label
and the bottleneck entries of \(\mathfrak A\).

For a structure involving several TMDs, the observable accuracy is bounded by
the least accurate required ingredient.  A linearly polarized gluon
coefficient known to higher order cannot upgrade an unknown process-specific
hard-color factor.

\section{Uncertainty propagation and convergence}
\label{sec:uncertainty}

\subsection{Shared pushforward}

Let \(\theta\) denote calibrated microscopic and matching parameters and let
\(\nu\) collect controlled discrete or continuous truncation choices.  The
prediction is
\begin{equation}
 \sigma_i=\sigma_i(\theta,\nu),
 \qquad
 \operatorname{Cov}(\sigma_i,\sigma_j)
 =\left\langle
 (\sigma_i-\bar\sigma_i)
 (\sigma_j-\bar\sigma_j)
 \right\rangle_{\theta,\nu}.
 \label{eq:shared-covariance}
\end{equation}
Member identity is retained through all transformations.  The covariance
therefore includes cross-flavor, quark--gluon, TMD--GPD, nucleon--deuteron,
rank, process, and scale correlations.

\subsection{Noninterchangeable uncertainty axes}

The following are reported separately before any optional combined summary:
\begin{enumerate}
\item microscopic parameter/posterior uncertainty;
\item Hamiltonian basis, Fock, and regulator truncation;
\item LF-to-QCD matching and missing-operator uncertainty;
\item perturbative coefficient and anomalous-dimension truncation;
\item evolution path/curl and scale variation;
\item nonperturbative CS-kernel uncertainty;
\item large-\(b\) boundary/discrepancy uncertainty;
\item collinear multiparton and threshold evolution uncertainty;
\item nuclear many-body and coherent-sector truncation;
\item process-factorization and Glauber status;
\item fixed-order and \(Y\)-term matching;
\item numerical transform, grid, interpolation, and quadrature error.
\end{enumerate}
Convergence in one axis cannot be used to hide nonconvergence in another.

\subsection{Observable-level convergence}

For a directed sequence of calculations \(n\), define
\begin{equation}
 \delta_i^{(n+1,n)}
 =
 \frac{
 \abs{\sigma_i^{(n+1)}-\sigma_i^{(n)}}
 }{
 \max(\abs{\sigma_i^{(n+1)}},\sigma_{i,\rm ref})
 }.
 \label{eq:observable-convergence}
\end{equation}
The sequence can increase Hamiltonian resolution, Fock content, matching order,
perturbative order, \(b\)-grid range, or fixed-order accuracy.  Each direction
has its own expected convergence law.  A smooth result at one resolution is
not a convergence study.

\section{Formal software architecture}
\label{sec:software}

\subsection{Core objects}

\begin{longtable}{L{0.27\textwidth}L{0.65\textwidth}}
\toprule
Object & Responsibility\\
\midrule
\endhead
\texttt{TMDSchemeId} & Complete tuple of UV scheme, rapidity regulator,
soft partition, staple/cusp convention, color representation, and tensor
conventions.\\
\texttt{MatchingBasis} & Closed LF and QCD operator bases, quantum numbers,
mixing blocks, local moments, and missing-operator status.\\
\texttt{LFtoQCDMatchingMap} & Versioned convolution/mixing kernel, matching
conditions, calibration states, step-scaling links, and discrepancy operators.\\
\texttt{RenormalizedTMDOperator} & Decorated operator identity with
\((\mu,\zeta)\), scheme, link/color, rank, target, nuclear sector, and member.\\
\texttt{AnomalousDimensionSet} & Cusp, noncusp, rapidity, beta-function, and
convention adapters with perturbative orders.\\
\texttt{EvolutionPath} & Endpoints, contour or \(\zeta\)-prescription,
profiles, finite-order curl treatment, and numerical solver metadata.\\
\texttt{CollinsSoperKernel} & Quark/gluon perturbative coefficients,
large-\(b\) completion, calibration members, representation status, and
matching profile.\\
\texttt{SmallBOPE} & Rank- and link-resolved coefficient matrices, collinear
operator basis, twist, lowest order, power remainder, and scheme.\\
\texttt{CollinearEvolution} & Nonsinglet, singlet, helicity, transversity,
double-flip, and twist-3 kernels plus threshold maps.\\
\texttt{RankTransform} & Forward/inverse Bessel transforms, phases, mass
powers, quadrature, analytic oracles, and round-trip reports.\\
\begin{tabular}[t]{@{}l@{}}\texttt{ProcessFactorization}\\\texttt{Record}\end{tabular} & Hard channel, TMD/FF basis, link/color
map, measurement function, Glauber and factorization status.\\
\texttt{WTerm} & Rank-resolved resummed contraction with scale and accuracy
manifest.\\
\texttt{YTerm} & Fixed-order result, asymptotic subtraction, common scheme,
kinematics, order, and profile.\\
\texttt{AccuracyManifest} & Full tuple \(\mathfrak A\), scale-variation plan,
bottleneck order, and generated human label.\\
\texttt{EvolutionEnsembleStore} & Common microscopic member identity and all
matching, kernel, scale, threshold, nuclear, and numerical axes.\\
\texttt{ClosureReport} & Marginal, moment, path, scheme, sum-rule, process,
\(W+Y\), and regression residuals.\\
\bottomrule
\end{longtable}

\subsection{Fail-closed composition}

A map can execute only if all required fields agree or an explicit adapter is
registered.  Mandatory comparisons include:
\begin{itemize}
\item \(\bD\) versus \(\bM\);
\item rank, Bessel order, Fourier phase, and reference mass;
\item quark versus gluon representation;
\item target and parton operator polarization;
\item future/past and ordered gluon link identity;
\item \(f/d\) color class;
\item UV, rapidity, and soft scheme;
\item \(\mu,\zeta,n_f\), and threshold history;
\item microscopic and nuclear member identity;
\item perturbative accuracy and fixed-order subtraction order;
\item process factorization and Glauber status.
\end{itemize}
An error reports the source and target types, missing adapter, and provenance
path.  It never substitutes a nominally similar function.

\subsection{Machine-readable manifests}

Every authoritative prediction is reproducible from four immutable manifests:
\begin{enumerate}
\item operator and matching manifest;
\item evolution and accuracy manifest;
\item process and \(W+Y\) manifest;
\item ensemble and numerical-convergence manifest.
\end{enumerate}
Their content hashes are embedded in the output metadata.  A plot without the
manifests is illustrative, not an authoritative artifact.


\section{Validation ladder and required property tests}
\label{sec:validation}

\subsection{Layered validation}

The matching and evolution layer is accepted in increasing dynamical depth:
\begin{enumerate}
\item \textbf{Type layer:} complete scheme, coordinate, rank, link/color,
      threshold, member, and process identity.
\item \textbf{Operator layer:} correct mixing blocks, local-current limits,
      Ward identities, and step-scaling composition.
\item \textbf{Evolution layer:} integrability, contour transitivity,
      representation dependence, and threshold closure.
\item \textbf{OPE layer:} small-\(b\) matching, collinear evolution
      compatibility, and rank-aware power counting.
\item \textbf{Parent layer:} common TMD/GPD/PDF/current reductions from the
      same matched microscopic operator.
\item \textbf{Nuclear layer:} evolution/composition commutation where valid,
      and separate evolution of many-body operators where not.
\item \textbf{Process layer:} hard/link/color qualification, \(W+Y\)
      fixed-order recovery, and Glauber gates.
\item \textbf{Predictive layer:} multi-scale withheld observables without
      per-TMD retuning.
\end{enumerate}

\subsection{Mandatory property tests}

At minimum, the implementation must test:
\begin{enumerate}
\item finite scheme transformations cancel in complete observables;
\item LF-to-QCD matching reproduces all calibration matrix elements and
      overconstrained holdout matrix elements within the error model;
\item step-scaling maps satisfy the cocycle relation;
\item local vector charge and total momentum are preserved;
\item the TMD evolution curl vanishes exactly in analytic benchmarks and is
      reported at finite perturbative order;
\item evolution is transitive and invertible within numerical tolerance;
\item future/past link-even equality and link-odd sign reversal are preserved;
\item quark and gluon kernels cannot be interchanged;
\item rank projectors commute with radial evolution;
\item every rank transform passes independent analytic forward/inverse tests;
\item small-\(b\) OPE followed by evolution agrees with evolution followed by
      coefficient matching;
\item unpolarized, helicity, transversity, and twist-3 kernels are not mixed;
\item singlet evolution preserves the momentum sum rule;
\item heavy-threshold matching preserves physical moments and process
      continuity to the retained order;
\item the spin-1 \(LL\) tree-level convention reproduces the helicity-state
      tensor difference with the declared sign adapter;
\item gluon \(f/d\) operator channels remain independent until a process map
      supplies hard weights;
\item the nuclear impulse convolution commutes with evolution under the
      theorem's assumptions;
\item an explicit two-body or pion operator fails the impulse-commutation test
      unless its own evolution block is supplied;
\item the asymptotic expansion of \(W\) reproduces all fixed-order singular
      terms through the declared order;
\item \(W+Y\) reproduces the fixed-order expansion and remains finite in the
      matching region;
\item process objects marked \texttt{BROKEN} or \texttt{UNASSESSED} cannot
      generate authoritative universal-TMD predictions;
\item all uncertainty axes preserve common member identity;
\item authoritative output hashes are unchanged by metadata-only migrations.
\end{enumerate}

\subsection{Deliberate negative injections}

The test suite must deliberately inject and reject:
\begin{itemize}
\item \(\bD\) used as \(\bM\);
\item a \(J_0\) transform applied to nonzero rank;
\item an inconsistent Fourier phase or reference mass;
\item a quark CS kernel applied to a gluon TMD;
\item a future-link Sivers boundary paired with a past-link process map without
      the sign adapter;
\item an \(f\)-type gluon coefficient combined with a \(d\)-type operator;
\item an \(n_f\) change without a threshold map;
\item a \(\MSbar\) coefficient combined with a different UV scheme;
\item a square-root soft TMD paired with a hard factor in an incompatible
      soft partition;
\item collinear and TMD evolution at inconsistent perturbative orders;
\item a twist-2 DGLAP kernel applied to a twist-3 multiparton vector;
\item an unrenormalized derivative moment promoted to a Qiu--Sterman function;
\item a polarization-dependent CS kernel introduced with no operator change;
\item a pilot or benchmark boundary promoted to production;
\item a \(Y\) subtraction built from different kinematics or a different
      fixed order;
\item a gluon-sensitive colored-final-state process with unassessed Glauber
      status;
\item a microscopic member mixed with a different nuclear member under one
      purported correlated sample.
\end{itemize}

\section{Analytic and numerical benchmark suite}
\label{sec:benchmarks}

\subsection{Benchmark A: exact two-scale potential}

Choose
\begin{equation}
 \cD(b;\mu)=a\ln\frac{\mu}{\mu_b}+d_0(b),
 \qquad
 \gamma_F(\mu,\zeta)=g_0-a\ln\frac{\sqrt\zeta}{\mu}.
 \label{eq:benchmark-a-anom}
\end{equation}
Then \(\Gamma=a\) and the evolution one-form is closed.  Evolve around a
rectangle in the scale plane and require exact return to the initial value.
Perturb one coefficient deliberately and verify that the loop integral equals
the area integral of the injected curl.

\subsection{Benchmark B: rank-\texorpdfstring{\(m\)}{m} Gaussian oracle}

For
\begin{equation}
 F^{(m)}(k)=k^m e^{-\alpha k^2},
 \label{eq:rank-gaussian}
\end{equation}
the Bessel transform is analytic and proportional to
\(b^m e^{-b^2/(4\alpha)}\).  Test ranks \(m=0,1,2,3\), all phase conventions,
and forward/inverse transforms on independent implementations.  Multiplying
by a radial evolution factor must preserve the harmonic.

\subsection{Benchmark C: singlet mixing with a conserved moment}

Construct a two-channel kernel matrix whose first moment has left null vector
\((1,1)\).  Evolve a quark--gluon vector and verify exact momentum
conservation.  Inject one incorrect off-diagonal sign and require the momentum
ledger to fail.

\subsection{Benchmark D: finite scheme transformation}

Select an invertible finite matrix \(\bm Z\) in a two-operator basis.  Transform
both TMDs and the hard factor according to
\cref{eq:finite-scheme-transform,eq:hard-scheme-transform}.  The cross section
must remain invariant.  Transform only the TMDs and verify a controlled
failure.

\subsection{Benchmark E: LF-to-QCD step scaling}

Use three finite-dimensional regulator levels with known embedding and
counterterm maps.  Construct \(\bm\Sigma_{2\leftarrow1}\),
\(\bm\Sigma_{3\leftarrow2}\), and \(\bm\Sigma_{3\leftarrow1}\), and test the
cocycle \cref{eq:step-scaling-cocycle}.  An omitted induced operator should
produce a nonzero cocycle defect and a failed holdout matrix element.

\subsection{Benchmark F: nuclear convolution commutation}

Choose analytic beta-function-like spectral and parton distributions so the
Mellin convolutions are known.  Verify
\begin{equation}
 P\otimesx(S\otimesx f)=S\otimesx(P\otimesx f).
 \label{eq:analytic-nuclear-associativity}
\end{equation}
Then introduce a scale-dependent two-body kernel and show that the naive
commutation fails unless its evolution block is included.

\subsection{Benchmark G: additive \texorpdfstring{\(W+Y\)}{W+Y} toy model}

Define an exactly known fixed-order spectrum and a resummed \(W\) term whose
expansion reproduces its singular part.  Construct \(Y\) with
\cref{eq:y-definition}.  Test small-\(q_T\) finiteness, fixed-order recovery,
profile independence through the retained order, and sensitivity to a
deliberately missing singular coefficient.

\subsection{Benchmark H: threshold continuity}

Use a toy heavy-flavor matching matrix and running coupling.  Evolve a
conserved moment through the threshold, verify physical continuity, reverse
the map, and test that changing \(n_f\) without the matching matrix fails.

\subsection{Benchmark I: link-reversal preservation}

Initialize
\begin{equation}
 F_{\even}^{[+]}=F_{\even}^{[-]},
 \qquad
 F_{\odd}^{[+]}=-F_{\odd}^{[-]}.
 \label{eq:benchmark-link-pair}
\end{equation}
Evolve both with the same representation kernel and verify the identities at
all scales.  Applying different kernels must be detected as an operator
identity mismatch, not accepted as physical link dependence.

\section{Staged implementation program}
\label{sec:stages}

The formal volume is independent of repository timing.  The following
packages are conceptual and should receive repository-specific Codex prompts
only after the preceding package reports its actual APIs and unresolved gaps.

\subsection{M0: matching and scheme skeleton}

\begin{enumerate}
\item Implement \texttt{TMDSchemeId}, \texttt{MatchingBasis}, and complete
      matching/operator manifests.
\item Encode the current accepted scheme and heterogeneous evolution status
      without changing numerical outputs.
\item Add finite scheme adapters and negative mismatch tests.
\item Keep all new microscopic matching disconnected from accepted production.
\end{enumerate}

\subsection{M1: analytic LF-to-QCD matching pilot}

\begin{enumerate}
\item Match the analytic common-overlap pilot to a regulated continuum operator
      basis.
\item Implement step scaling and local-current constraints.
\item Demonstrate overconstrained matching on quark, antiquark, and gluon
      benchmark states.
\item Retain explicit missing-operator and truncation residuals.
\end{enumerate}

\subsection{M2: two-scale evolution engine}

\begin{enumerate}
\item Implement anomalous-dimension sets and convention adapters.
\item Implement contour, integrability-improved, and optional
      \(\zeta\)-prescription paths.
\item Report finite-order curl and transitivity residuals.
\item Add separate quark and gluon kernel objects and large-\(b\) interfaces.
\end{enumerate}

\subsection{M3: rank-aware small-\texorpdfstring{\(b\)}{b} OPE and collinear evolution}

\begin{enumerate}
\item Encode the twist-2 and twist-3 operator bases and coefficient status.
\item Implement nonsinglet, singlet, helicity, transversity, double-flip, and
      multiparton evolution interfaces.
\item Add threshold maps and sum-rule tests.
\item Validate rank-aware transforms and derivative-moment renormalization.
\end{enumerate}

\subsection{M4: matched nuclear evolution}

\begin{enumerate}
\item Apply evolution to the resolved nucleon/deuteron operator graph while
      retaining all component and member labels.
\item Prove numerical commutation for impulse terms.
\item Implement separate evolution blocks for pion, coherent, and two-body
      operators as they become available.
\item Reject unresolved scale dependence hidden in a scalar nuclear response.
\end{enumerate}

\subsection{M5: process records and \texorpdfstring{\(W+Y\)}{W+Y}}

\begin{enumerate}
\item Implement qualified Drell--Yan and SIDIS records first.
\item Add inclusive \(b_1\) and local-current closure.
\item Implement rank-resolved fixed-order asymptotic subtraction.
\item Add gluon-sensitive channels only with explicit factorization and
      Glauber status.
\item Reserve at least one multi-scale observable family for withheld
      validation.
\end{enumerate}

\section{Risks, failure modes, and explicit nonclaims}
\label{sec:risks}

\begin{longtable}{L{0.22\textwidth}L{0.34\textwidth}L{0.36\textwidth}}
\toprule
Risk & Failure mode & Required safeguard\\
\midrule
\endhead
Matching by name & Assigning one finite factor to each TMD without an operator
basis. & Match closed operator blocks and overconstrain shared parameters with
multiple matrix elements.\\
Scheme drift & Combining soft, rapidity, hard, or Fourier conventions from
different formulations. & Complete scheme tuple and finite adapters; fail
closed on incomplete identity.\\
Factor-of-two evolution error & Mixing derivatives in \(\ln\zeta\),
\(\ln\sqrt\zeta\), \(\ln\mu\), and \(\ln\mu^2\). & Store derivative convention
and verify the mixed-derivative identity analytically.\\
Decorative high accuracy & Quoting \(\mathrm{N}^k\mathrm{LL}\) from the cusp
order while other ingredients are lower. & Full accuracy tuple and bottleneck
report.\\
Large-\(b\) overfitting & Giving every TMD an independent CS kernel or width. &
Representation-level kernel plus shared microscopic boundary and withheld
cross-channel tests.\\
Casimir overreach & Imposing quadratic Casimir scaling beyond its verified
order or on the nonperturbative kernel. & Independent quark/gluon records and
order-qualified generalized relations.\\
Naive marginal equality & Treating \(\int\dd^2k_T\) of a subtracted TMD as a
PDF/GPD/current. & Explicit small-\(b\) OPE and matched commuting squares.\\
Twist confusion & Evolving Sivers or Boer--Mulders with a twist-2 DGLAP kernel.
& Typed collinear operator basis and multiparton twist-3 evolution.\\
Rank loss & Applying \(J_0\) or scalar asymptotic subtraction to all
modulations. & Rank identity on every TMD, hard tensor, transform, and \(Y\)
term.\\
Nuclear reassembly & Evolving an already summed deuteron curve and losing
constituent/mechanism correlations. & Evolve the resolved component graph and
preserve member identity.\\
Soft double counting & Including a microscopic eikonal exchange in both the
boundary and CS/soft sector. & Phase budget and explicit overlap subtraction.\\
Threshold discontinuity & Changing \(n_f\) without coefficient, kernel, and
coupling matching. & Versioned threshold map and conserved-moment tests.\\
Missing \(Y\) physics & Retuning the nonperturbative TMD to fit moderate/high
\(q_T\). & Independent fixed-order calculation and same-order asymptotic
subtraction.\\
Factorization overreach & Applying universal TMDs in a color-entangled or
Glauber-obstructed process. & Process status gate and generalized soft matrix
where justified.\\
Positivity clipping & Repairing finite-order scheme negativity by projecting
onto a positive cone. & Diagnose the layer; preserve the unmodified matched
operator and test physical observables.\\
\bottomrule
\end{longtable}

The volume does not claim:
\begin{itemize}
\item that the LF-to-QCD matching is perturbative over the full hadronic
      domain;
\item that one nonperturbative CS kernel has already been determined for both
      quarks and gluons;
\item that every spin-1 small-\(b\) coefficient is presently known;
\item that twist-3 multiparton evolution can be reduced exactly to one
      diagonal function;
\item that standard TMD factorization holds for every gluon-sensitive final
      state;
\item that the accepted fixed-scale phenomenological boundary already obeys
      one common evolution scheme;
\item that successful evolution alone converts that boundary into a
      microscopic prediction.
\end{itemize}

\section{Acceptance criteria and handoff}
\label{sec:acceptance}

\subsection{Formal acceptance criteria}

The Volume~V layer is complete only when:
\begin{enumerate}
\item every regulated parent operator has a complete LF and QCD identity;
\item a shared LF-to-QCD matching map is determined or explicitly unavailable
      for each required operator block;
\item charges, Ward identities, and momentum/EMT moments close after matching;
\item step-scaling and regulator convergence are demonstrated;
\item the TMD scheme, soft partition, rapidity regulator, and scale convention
      are immutable metadata;
\item the two-scale evolution field satisfies exact analytic integrability and
      reports finite-order curl defects;
\item quark and gluon CS kernels are separate and target-polarization
      independent at the declared leading power;
\item all transverse ranks use correct Bessel kernels, phases, and mass powers;
\item every small-\(b\) coefficient identifies its collinear operator, twist,
      link/color class, lowest known order, and power remainder;
\item nonsinglet, singlet, helicity, transversity, double-flip, twist-3, and
      threshold evolution routes pass their sum-rule tests;
\item the resolved microscopic and nuclear member identity survives matching
      and evolution;
\item evolution and impulse nuclear composition commute where the theorem
      applies, and noncommuting mechanisms have their own operator evolution;
\item process records carry hard, link, color, soft, measurement, Glauber, and
      factorization information;
\item rank-resolved \(W+Y\) reproduces fixed order through the declared order;
\item scheme and scale variations do not erase or reshuffle provenance;
\item at least one multi-scale nucleon family and one deuteron family withheld
      from matching/calibration are successfully predicted;
\item all structural, marginal, sum-rule, process, numerical, and immutable
      regression gates pass.
\end{enumerate}

\subsection{Contract to the inference volume}

The next formal layer receives:
\begin{itemize}
\item a common microscopic ensemble;
\item matching and evolution manifests;
\item rank-, flavor-, link/color-, and mechanism-resolved predictions at all
      requested scales;
\item a separated uncertainty budget;
\item process-qualified observable likelihoods;
\item withheld-data designations fixed before final calibration.
\end{itemize}
It must infer shared microscopic, matching, and controlled discrepancy
parameters without rebuilding one latent normalization per TMD.

\subsection{Compact adoption questions}

Before accepting an evolved observable, ask:
\begin{enumerate}
\item Which regulated operator produced it?
\item Which matching map and calibration matrix elements define it?
\item Which complete QCD scheme and scales does it inhabit?
\item Which rank and Fourier convention are used?
\item Which collinear operator and OPE coefficient control small \(b\)?
\item Which evolution family and threshold history were applied?
\item Does the same microscopic member survive the full chain?
\item Which process theorem supplies the link and hard/color map?
\item What fixed-order calculation and asymptotic subtraction define \(Y\)?
\item Which uncertainty axis dominates, and is it distinct from numerical
      error?
\item Does a withheld observable test the same shared parameters?
\end{enumerate}
Any unanswered question blocks an authoritative claim.

\appendix

\section{Convention dictionary for two-scale evolution}
\label{app:conventions}

The project equations are
\begin{equation}
 \frac{\dd\ln F}{\dd\ln\mu}=\gamma_F,
 \qquad
 \frac{\dd\ln F}{\dd\ln\sqrt\zeta}=-\cD,
 \qquad
 \frac{\dd\cD}{\dd\ln\mu}=\Gamma.
 \label{eq:project-evolution-convention}
\end{equation}
Common alternative notations include
\begin{equation}
 \widetilde K=-\cD,
 \qquad
 \gamma_K=2\Gamma,
 \qquad
 \frac{\dd}{\dd\ln\mu^2}=\frac12\frac{\dd}{\dd\ln\mu},
 \qquad
 \frac{\dd}{\dd\ln\zeta}=\frac12
 \frac{\dd}{\dd\ln\sqrt\zeta}.
 \label{eq:alternative-evolution-notation}
\end{equation}
Every adapter is tested by transforming all four relations together.  A
single symbol substitution is insufficient.

\section{Mellin-convolution identities}
\label{app:mellin}

For
\begin{equation}
 [A\otimesx B](x)
 =\int_x^1\frac{\dd z}{z}A(z)B(x/z),
 \label{eq:mellin-definition}
\end{equation}
associativity follows under standard integrability assumptions:
\begin{equation}
 A\otimesx(B\otimesx C)
 =(A\otimesx B)\otimesx C.
 \label{eq:mellin-associativity}
\end{equation}
This identity underlies the evolution--impulse commutation theorem.  Matrix
convolutions preserve the order of operator indices even though the scalar
Mellin integration is associative.

\section{Minimal accuracy-manifest schema}
\label{app:accuracy-schema}

An authoritative manifest contains at least:
\begin{verbatim}
{
  "scheme_id": "...",
  "parton_representation": "FUNDAMENTAL|ADJOINT",
  "cusp_order": 0,
  "noncusp_order": 0,
  "rapidity_order": 0,
  "beta_order": 0,
  "small_b_matching_order": 0,
  "hard_order": 0,
  "splitting_order": 0,
  "fixed_order_cross_section": 0,
  "y_subtraction_order": 0,
  "evolution_prescription": "...",
  "finite_order_curl_treatment": "...",
  "profile_id": "...",
  "threshold_history": [],
  "scale_variation_plan": "...",
  "bottleneck": "..."
}
\end{verbatim}
The human-readable logarithmic label is derived from this object.

\section{Minimal process-factorization schema}
\label{app:process-schema}

\begin{verbatim}
{
  "process_id": "...",
  "external_states": [],
  "hard_channel": "...",
  "measurement": "...",
  "tmd_operator_ids": [],
  "fragmentation_operator_ids": [],
  "link_map": {},
  "gluon_color_weights": {},
  "soft_scheme": "...",
  "scale_relation": "...",
  "power_counting": "...",
  "glauber_status": "...",
  "factorization_status": "...",
  "fixed_order_reference": "...",
  "asymptotic_subtraction": "...",
  "rank_harmonics": [],
  "accuracy_manifest_hash": "..."
}
\end{verbatim}

\section{Summary of stable formal requirements}
\label{app:requirements}

\begin{longtable}{L{0.16\textwidth}L{0.76\textwidth}}
\toprule
ID & Requirement\\
\midrule
\endhead
V5.TYPE.1 & Keep \(\bD\), \(\bM\), \(\qT\), and loop momenta as distinct types.\\
V5.TYPE.2 & Make rank, Bessel order, Fourier phase, and mass convention part of operator identity.\\
V5.SCHEME.1 & Store a complete UV/rapidity/soft/link/color scheme tuple.\\
V5.SCHEME.2 & Transform TMDs, hard factors, and \(Y\) subtractions covariantly.\\
V5.MATCH.1 & Match closed operator bases, not named curves.\\
V5.MATCH.2 & Overconstrain shared matching parameters with multiple matrix elements.\\
V5.MATCH.3 & Preserve local currents, Ward identities, and EMT moments.\\
V5.MATCH.4 & Demonstrate regulator step scaling and cocycle closure.\\
V5.EVOL.1 & Use convention-consistent \((\mu,\zeta)\) equations and mixed derivatives.\\
V5.EVOL.2 & Report finite-order curl, contour, and transitivity residuals.\\
V5.EVOL.3 & Keep quark and gluon CS kernels separate; do not fit one per polarization.\\
V5.RANK.1 & Use \(J_{|m|}\) and declared rank normalizations for every transform and process harmonic.\\
V5.OPE.1 & Identify collinear operator, twist, link/color class, coefficient order, and power remainder.\\
V5.OPE.2 & Distinguish twist-2 from multiparton twist-3 evolution.\\
V5.COLL.1 & Separate nonsinglet, singlet, helicity, transversity, double-flip, and twist-3 kernels.\\
V5.COLL.2 & Use explicit heavy-flavor threshold maps and conserve physical moments.\\
V5.NUC.1 & Evolve the resolved nuclear graph and prove impulse commutation where valid.\\
V5.NUC.2 & Give many-body/coherent operators their own matching and evolution.\\
V5.PROC.1 & Require hard, link, color, soft, measurement, Glauber, and factorization records.\\
V5.WY.1 & Define \(Y\) from the same-order, same-scheme asymptotic expansion of \(W\).\\
V5.WY.2 & Perform matching separately for each rank and angular harmonic.\\
V5.UNC.1 & Preserve microscopic member identity and separated uncertainty axes.\\
V5.ACC.1 & Reserve multi-scale nucleon and deuteron families for withheld validation.\\
\bottomrule
\end{longtable}

\small
\begin{thebibliography}{99}

\bibitem{CollinsSoper1981}
J.~C.~Collins and D.~E.~Soper,
``Back-to-back jets in QCD,''
Nucl.\ Phys.\ B \textbf{193}, 381 (1981);
Erratum Nucl.\ Phys.\ B \textbf{213}, 545 (1983).

\bibitem{CollinsSoper1982}
J.~C.~Collins and D.~E.~Soper,
``Parton distribution and decay functions,''
Nucl.\ Phys.\ B \textbf{194}, 445 (1982).

\bibitem{CSS1985}
J.~C.~Collins, D.~E.~Soper, and G.~F.~Sterman,
``Transverse momentum distribution in Drell--Yan pair and W and Z boson
production,''
Nucl.\ Phys.\ B \textbf{250}, 199 (1985).

\bibitem{Collins2011}
J.~C.~Collins,
\emph{Foundations of Perturbative QCD},
Cambridge University Press (2011).

\bibitem{AybatRogers2011}
S.~M.~Aybat and T.~C.~Rogers,
``TMD parton distribution and fragmentation functions with QCD evolution,''
Phys.\ Rev.\ D \textbf{83}, 114042 (2011), arXiv:1101.5057.

\bibitem{EIS2012}
M.~G.~Echevarria, A.~Idilbi, and I.~Scimemi,
``Factorization theorem for Drell--Yan at low \(q_T\) and transverse-momentum
distributions on the light cone,''
JHEP \textbf{07}, 002 (2012), arXiv:1111.4996.

\bibitem{EIS2013}
M.~G.~Echevarria, A.~Idilbi, and I.~Scimemi,
``Soft and collinear factorization and transverse momentum dependent parton
distribution functions,''
Phys.\ Lett.\ B \textbf{726}, 795 (2013), arXiv:1211.1947.

\bibitem{Chiu2012}
J.-Y.~Chiu, A.~Jain, D.~Neill, and I.~Z.~Rothstein,
``A formalism for the systematic treatment of rapidity logarithms in quantum
field theory,''
JHEP \textbf{05}, 084 (2012), arXiv:1202.0814.

\bibitem{EchevarriaScimemiVladimirov2016}
M.~G.~Echevarria, I.~Scimemi, and A.~Vladimirov,
``Universal transverse momentum dependent soft function at NNLO,''
Phys.\ Rev.\ D \textbf{93}, 054004 (2016), arXiv:1511.05590.

\bibitem{EchevarriaGTMD2016}
M.~G.~Echevarria, A.~Idilbi, K.~Kanazawa, C.~Lorc\'e, A.~Metz,
B.~Pasquini, and M.~Schlegel,
``Proper definition and evolution of generalized transverse momentum
distributions,''
Phys.\ Lett.\ B \textbf{759}, 336 (2016), arXiv:1602.06953.

\bibitem{ScimemiVladimirov2017}
I.~Scimemi and A.~Vladimirov,
``Analysis of vector boson production within TMD factorization,''
Eur.\ Phys.\ J.\ C \textbf{78}, 89 (2018), arXiv:1706.01473.

\bibitem{ScimemiVladimirov2018}
I.~Scimemi and A.~Vladimirov,
``Systematic analysis of double-scale evolution,''
JHEP \textbf{08}, 003 (2018), arXiv:1803.11089.

\bibitem{GutierrezReyes2017}
D.~Gutierrez-Reyes, I.~Scimemi, and A.~A.~Vladimirov,
``Twist-2 matching of transverse momentum dependent distributions,''
Phys.\ Lett.\ B \textbf{769}, 84 (2017), arXiv:1702.06558.

\bibitem{ScimemiVladimirovTwist3}
I.~Scimemi and A.~Vladimirov,
``Matching of transverse momentum dependent distributions at twist-3,''
Eur.\ Phys.\ J.\ C \textbf{78}, 802 (2018), arXiv:1804.08148.

\bibitem{VladimirovMoosScimemi2021}
A.~Vladimirov, V.~Moos, and I.~Scimemi,
``Transverse momentum dependent operator expansion at next-to-leading power,''
JHEP \textbf{01}, 110 (2022), arXiv:2109.09771.

\bibitem{GutierrezReyes2019}
D.~Gutierrez-Reyes, S.~Leal-Gomez, I.~Scimemi, and A.~Vladimirov,
``Linearly polarized gluons at next-to-next-to-leading order and the Higgs
transverse momentum distribution,''
JHEP \textbf{11}, 121 (2019), arXiv:1907.03780.

\bibitem{LuoYangZhuZhu2019}
M.-X.~Luo, T.-Z.~Yang, H.~X.~Zhu, and Y.~J.~Zhu,
``Transverse parton distribution and fragmentation functions at NNLO: the
gluon case,''
JHEP \textbf{01}, 040 (2020), arXiv:1909.13820.

\bibitem{MoultZhuZhu2022}
I.~Moult, H.~X.~Zhu, and Y.~J.~Zhu,
``The four-loop QCD rapidity anomalous dimension,''
arXiv:2205.02249.

\bibitem{DuhrMistlbergerVita2022}
C.~Duhr, B.~Mistlberger, and G.~Vita,
``The four-loop rapidity anomalous dimension and event shapes to fourth
logarithmic order,''
arXiv:2205.02242.

\bibitem{ZhuLinearlyPolarized2025}
Y.~J.~Zhu,
``The \(\mathrm{N}^3\mathrm{LO}\) twist-2 matching of linearly polarized gluon
TMDs,''
arXiv:2509.01703 (2025).

\bibitem{ZhuTransversity2025}
Y.~J.~Zhu,
``The \(\mathrm{N}^3\mathrm{LO}\) twist-2 matching of TMD quark transversity,''
arXiv:2509.17568 (2025).

\bibitem{ZhuDGLAP2026}
Y.~J.~Zhu,
``NNLO DGLAP splitting functions from collinear matching of TMDs,''
arXiv:2603.04039 (2026).

\bibitem{KangXiaoYuan2011}
Z.-B.~Kang, B.-W.~Xiao, and F.~Yuan,
``QCD resummation for single spin asymmetries,''
Phys.\ Rev.\ Lett.\ \textbf{107}, 152002 (2011), arXiv:1106.0266.

\bibitem{SunYuan2013}
P.~Sun and F.~Yuan,
``TMD evolution: matching SIDIS to Drell--Yan and W/Z boson production,''
Phys.\ Rev.\ D \textbf{88}, 114012 (2013), arXiv:1308.5003.

\bibitem{KorchemskySterman1995}
G.~P.~Korchemsky and G.~F.~Sterman,
``Nonperturbative corrections in resummed cross sections,''
Nucl.\ Phys.\ B \textbf{437}, 415 (1995), arXiv:hep-ph/9411211.

\bibitem{ScimemiVladimirovRenormalon2017}
I.~Scimemi and A.~Vladimirov,
``Power corrections and renormalons in transverse momentum distributions,''
JHEP \textbf{03}, 002 (2017), arXiv:1609.06047.

\bibitem{Avkhadiev2024}
A.~Avkhadiev, P.~E.~Shanahan, M.~L.~Wagman, and Y.~Zhao,
``Determination of the Collins--Soper kernel from lattice QCD,''
arXiv:2402.06725 (2024).

\bibitem{Bollweg2024}
D.~Bollweg, X.~Gao, S.~Mukherjee, and Y.~Zhao,
``Nonperturbative Collins--Soper kernel from chiral quarks with physical
masses,''
arXiv:2403.00664 (2024).

\bibitem{TanEtAl2025}
J.-X.~Tan et al.,
``Lattice QCD determination of the Collins--Soper kernel in the continuum and
physical-mass limits,''
arXiv:2511.22547 (2025).

\bibitem{BacchettaMulders2000}
A.~Bacchetta and P.~J.~Mulders,
``Deep inelastic leptoproduction of spin-one hadrons,''
Phys.\ Rev.\ D \textbf{62}, 114004 (2000), arXiv:hep-ph/0007120.

\bibitem{KumanoSong2021}
S.~Kumano and Q.-T.~Song,
``Transverse-momentum-dependent parton distribution functions for spin-1
hadrons,''
Phys.\ Rev.\ D \textbf{103}, 014025 (2021), arXiv:2011.08583.

\bibitem{Boer2016}
D.~Boer, S.~Cotogno, T.~van Daal, P.~J.~Mulders, A.~Signori, and Y.-J.~Zhou,
``Gluon and Wilson-loop TMDs for hadrons of spin \(\leq1\),''
JHEP \textbf{10}, 013 (2016), arXiv:1607.01654.

\bibitem{HoodbhoyJaffeManohar1989}
P.~Hoodbhoy, R.~L.~Jaffe, and A.~Manohar,
``Novel effects in deep-inelastic scattering from spin-one hadrons,''
Nucl.\ Phys.\ B \textbf{312}, 571 (1989).

\end{thebibliography}

\end{document}
